\chapter{Fourier at the Crossroads}

\begin{center}\itshape
Profound study of nature is the most fertile source of mathematical discoveries.\\
\upshape --- Joseph Fourier, \emph{The Analytical Theory of Heat} (1822), Preliminary Discourse, in Alexander Freeman's translation (1878)
\end{center}

In 1785 the Marquis de Condorcet, mathematician, economist and soon a revolutionary, published a long \emph{Essai sur l'application
de l'analyse à la probabilité des décisions rendues à la pluralité des voix}. He wanted to know when a vote by majority is likely to
reach the right decision, and on the way he found something unsettling. Suppose three voters rank three options $a$, $b$, $c$ as
$a>b>c$, $b>c>a$ and $c>a>b$. Then a majority prefers $a$ to $b$, a majority prefers $b$ to $c$, and a majority prefers $c$ to $a$.
Each voter is perfectly rational, and the group goes round in a circle. This is the \emph{Condorcet paradox}.

For a century and a half it stayed a curiosity. Then in 1951 Kenneth Arrow, in his thesis \emph{Social Choice and Individual
Values}, proved that the circle cannot be avoided by being clever: with three or more options, every method of turning individual
rankings into a collective ranking that respects a few innocent-looking conditions is a dictatorship, the ranking of a single
voter. Arrow received the Nobel prize in economics in 1972, and his theorem became the founding result of social choice theory.

How often does the paradox actually bite? In 1952 Georges-Théodule Guilbaud, in a long paper in the journal \emph{Économie
appliquée}, gave the answer for three options and voters who rank them at random: as the number of voters grows, the probability
that majority rule produces a cycle tends to about $8.77\%$. He gave the formula with the remark that it follows ``by the usual
methods'', and no further explanation.

Fifty years later Gil Kalai looked at the problem with a new pair of glasses. A voting rule on two options is a function
$f:\{-1,1\}^n\to\{-1,1\}$, and every such function has a Fourier expansion, like a periodic signal. Kalai showed in 2002 that the
probability of a rational outcome is an explicit weighted sum of the squares of the Fourier coefficients of $f$. Guilbaud's number
falls out in two lines, and Arrow's theorem, in the case of three options, becomes the statement that only one kind of Fourier
spectrum is perfectly balanced. You will rebuild both in Problem~\ref{pb:arrow}.

This is the theme of the chapter: one idea, \emph{expand in characters}, working at every crossroads of the series. Characters of
$\Z/p\Z$ count solutions of equations and arithmetic progressions; characters of the circle measure how evenly a sequence is
spread and decide whether a random variable has a density; characters of the cube $\{-1,1\}^n$ analyse votes, cryptographic
functions and card shuffles; and the classical Fourier series of Tome~4 answer old extremal questions about trigonometric
polynomials, polygons and lattice points.

\begin{think}
Take the three voters above and let each one instead choose a ranking of $a$, $b$, $c$ uniformly at random among the six possible,
independently. What is the probability that majority rule produces a cycle? (There are $6^3=216$ profiles; count the bad ones.)
Then guess: as the number of voters grows, does the probability of a cycle go to $0$, to $1$, or somewhere in between?
\end{think}

Throughout the chapter we write $e(x)=e^{2\pi ix}$, so that $e(x+1)=e(x)$, and for a prime $p$ we identify $\mathbb F_p$ with
$\Z/p\Z$; the expression $e(\frac ap)$ then makes sense for $a\in\mathbb F_p$. The basic fact behind everything is orthogonality:
\[
\sum_{r\in\mathbb F_p}e\Bigl(\frac{ra}p\Bigr)=\begin{cases}p&a=0,\\0&a\ne0,\end{cases}
\]
so that $\frac1p\sum_re(\frac{r(X-x)}p)$ is the indicator of the equation $X=x$. Counting solutions becomes summing exponentials.

\section{Characters modulo $p$}

For a function $g:\mathbb F_p\to\C$ we write $\hat g(r)=\sum_xg(x)e(\frac{rx}p)$, and for a set $A$ we write $\hat A(r)$ for $\hat 1_A(r)$. Parseval's
identity (Tome~4, Theorem~7.3) reads $\sum_r|\hat g(r)|^2=p\sum_x|g(x)|^2$; for a set, $\sum_r|\hat A(r)|^2=p|A|$. The problems below use
little more than this, Cauchy--Schwarz, and the value of the Gauss sum.

\begin{exercise}[$\star\star$]\label{pb:bilinear}
\emph{Sums over boxes.} Let $p$ be an odd prime.
\begin{enumerate}[label=(\alph*)]
\item Let $A,B\subseteq\mathbb F_p$ and $r\ne0$. Show that
\[
\Bigl|\sum_{a\in A}\sum_{b\in B}e\Bigl(\frac{rab}p\Bigr)\Bigr|\le\sqrt{p|A||B|}.
\]
\item Let $U,V\subseteq\mathbb F_p^2$, let $x\ne0$, and let $N(x)$ be the number of pairs $(u,v)\in U\times V$ with $u\cdot v=x$ (the dot
product $u_1v_1+u_2v_2$). Show that
\[
\Bigl|N(x)-\frac{|U||V|}p\Bigr|\le\sqrt{p|U||V|}.
\]
\item Let $A,B,C,D\subseteq\mathbb F_p$ with $|A||B||C||D|>p^3$. Show that every $x\ne0$ can be written $ab+cd$ with $a\in A$, $b\in B$,
$c\in C$, $d\in D$. In particular, if $|A|>p^{3/4}$, every non-zero residue is of the form $ab+cd$ with $a,b,c,d\in A$.
\item Let $A=B=\{1,\dots,m\}$ with $m\le\frac{\sqrt p}3$. Show that the sum in (a) with $r=1$ has real part at least $\frac{m^2}2$. Taking
$m=\lfloor\frac{\sqrt p}3\rfloor$, conclude that (a) cannot be improved by more than a constant factor.
\end{enumerate}
\end{exercise}

\begin{exercise}[$\star\star\star$]\label{pb:kloosterman}
\emph{Kloosterman sums.} Let $p$ be an odd prime. For $a,b\in\mathbb F_p$ let
\[
K(a,b)=\sum_{x\in\mathbb F_p^\times}e\Bigl(\frac{ax+b\bar x}p\Bigr),
\]
where $\bar x$ is the inverse of $x$ modulo $p$, and $K(a)=K(a,1)$.
\begin{enumerate}[label=(\alph*)]
\item Show that $K(a)$ is real, that $K(0)=-1$, and that $K(a,b)=K(ab)$ when $b\ne0$.
\item Show that $\sum_{a\in\mathbb F_p}K(a)^2=p^2-p$.
\item Show that $\sum_{a\in\mathbb F_p}K(a)^4=2p^3-3p^2-3p$.
\item Deduce that $|K(a)|<1.19\,p^{3/4}$ for every $a$. How does this compare with the trivial bound, and with what (b) gives?
\end{enumerate}
\end{exercise}

\begin{exercise}[$\star\star$]\label{pb:apcount}
\emph{Counting progressions with Fourier.} Let $p$ be an odd prime, $A\subseteq\mathbb F_p$, and let $T(A)$ be the number of pairs
$(x,d)\in\mathbb F_p^2$ with $x$, $x+d$, $x+2d$ all in $A$ (the pairs with $d=0$ are included).
\begin{enumerate}[label=(\alph*)]
\item Show that $T(A)=\frac1p\sum_{r\in\mathbb F_p}\hat A(r)^2\hat A(-2r)$.
\item Deduce that $\bigl|T(A)-\frac{|A|^3}p\bigr|\le|A|\max_{r\ne0}|\hat A(r)|$.
\item Let $Q$ be the set of non-zero squares modulo $p$. Show that $\bigl|T(Q)-\frac{(p-1)^3}{8p}\bigr|\le\frac{(\sqrt p+1)(p-1)}4$; so $T(Q)=\frac{p^2}8+O(p^{3/2})$.
\item Let $|A|=\delta p$. Show that if $|\hat A(r)|\le\frac12\delta^2p$ for all $r\ne0$ and $\delta^2p>2$, then $A$ contains a
progression $x,x+d,x+2d$ with $d\ne0$.
\end{enumerate}
\end{exercise}

\begin{exercise}[$\star\star\star$]\label{pb:dftspectrum}
\emph{The eigenvalues of the discrete Fourier transform.} Let $n\ge2$ and let $F$ be the unitary DFT on $\C^n$:
$(Fx)_j=\frac1{\sqrt n}\sum_{k=0}^{n-1}x_ke(-\frac{jk}n)$, indices modulo $n$.
\begin{enumerate}[label=(\alph*)]
\item Show that $F^2$ is the reversal $x_k\mapsto x_{-k}$, that $F^4=I$, and that $F$ is diagonalisable with eigenvalues among
$1,-1,-i,i$. Call their multiplicities $m_1$, $m_{-1}$, $m_{-i}$, $m_i$.
\item Express $\operatorname{tr}F$ and $\operatorname{tr}F^2$ in terms of the multiplicities, and compute both traces.
\item Deduce the multiplicities: writing $n=4m+s$ with $0\le s\le3$, the vector $(m_1,m_{-1},m_{-i},m_i)$ is
\[
(m+1,m,m,m-1),\ (m+1,m,m,m),\ (m+1,m+1,m,m),\ (m+1,m+1,m+1,m)
\]
for $s=0,1,2,3$.
\end{enumerate}
\end{exercise}

\begin{exercise}[$\star\star$]\label{pb:newman}
\emph{Newman's phenomenon.} Let $s(n)$ be the sum of the binary digits of $n$, and $\omega=e(\frac13)$.
\begin{enumerate}[label=(\alph*)]
\item For $k\ge1$ and $j\in\{0,1,2\}$, show that the number of $n<2^k$ with $s(n)\equiv j\pmod3$ is $\frac13\bigl(2^k+2\cos\frac{\pi(k-2j)}3\bigr)$.
\item Show that
\[
\sum_{\substack{0\le n<4^k\\3\mid n}}(-1)^{s(n)}=2\cdot3^{k-1}\qquad\text{and}\qquad\sum_{\substack{0\le n<2\cdot4^k\\3\mid n}}(-1)^{s(n)}=3^k.
\]
So among the multiples of $3$, those with an even digit sum win, by a margin of order $N^{\log3/\log4}$ below $N$.
\end{enumerate}
\end{exercise}

\section{Equidistribution on the circle}

A sequence $x_1,x_2,\dots$ in $[0,1[$ is equidistributed when every interval gets its fair share of points; Weyl's criterion
(Tome~4, Theorem~6.13) says this happens exactly when the exponential sums $\frac1N\sum_ne(hx_n)$ tend to $0$ for every $h\ne0$. The
same sums measure random variables: the characteristic function of $X$ at $2\pi h$ is $\mathbb E[e(hX)]$. The problems of this
section make the criterion quantitative, apply it to random walks and lacunary sums, and use characters to decide when a
law has a density and when a sum equals an integral.

\begin{exercise}[$\star\star\star$]\label{pb:erdosturan}
\emph{An Erdős--Turán inequality.} Let $x_1,\dots,x_N\in[0,1[$, $S(h)=\sum_{n=1}^Ne(hx_n)$, and let $D$ be the \emph{discrepancy}, the
supremum over intervals $I\subseteq[0,1[$ of $\bigl|\frac1N\#\{n:x_n\in I\}-|I|\bigr|$.
\begin{enumerate}[label=(\alph*)]
\item Let $0<\delta<1$ and let $J$ be an interval of length $L\le1-\delta$. The $1$-periodic \emph{trapezoid} $t$ is the convolution
of the periodic indicator of $J$ with $\frac1\delta1_{[0,\delta]}$. Show that $0\le t\le1$, $\int_0^1t=L$, and that its Fourier
coefficients satisfy $|\hat t(h)|\le\min\bigl(\frac1{\pi|h|},\frac1{\pi^2h^2\delta}\bigr)$ for $h\ne0$.
\item Deduce that
\[
D\le2\delta+\frac2N\sum_{h\ge1}\min\Bigl(\frac1{\pi h},\frac1{\pi^2h^2\delta}\Bigr)|S(h)|.
\]
\item Let $p$ be an odd prime and $x_k=\{\frac{k^2}p\}$ for $0\le k<p$. Show that the discrepancy of these $p$ points is at most
$\frac{\log p+10}{\pi\sqrt p}$.
\end{enumerate}
\end{exercise}

\begin{exercise}[$\star\star$]\label{pb:circlewalk}
\emph{A random walk on the circle.} Let $\alpha$ be real, $\varepsilon_1,\varepsilon_2,\dots$ independent fair signs, and
$X_n=\{\alpha(\varepsilon_1+\cdots+\varepsilon_n)\}\in[0,1[$.
\begin{enumerate}[label=(\alph*)]
\item Show that if $\alpha$ is irrational, $X_n$ converges in law to the uniform distribution: $\mathbb E[g(X_n)]\to\int_0^1g$ for every
continuous $1$-periodic $g$.
\item Let $\alpha=\frac aq$ in lowest terms, $q\ge2$. Show that if $q$ is odd, $X_n$ converges in law to the uniform distribution on
$\{0,\frac1q,\dots,\frac{q-1}q\}$, and that if $q$ is even, $X_{2n}$ and $X_{2n+1}$ converge in law to the uniform distributions on two
disjoint sets of $\frac q2$ points, so that $X_n$ does not converge.
\end{enumerate}
\end{exercise}

\begin{exercise}[$\star\star\star$]\label{pb:kac}
\emph{Lacunary sums look Gaussian.} Let $t$ be uniform on $[0,2\pi]$ and, for $N\ge1$,
\[
S_N=\sum_{k=1}^N\cos(2^kt),\qquad T_N=\sum_{k=1}^N\cos(kt).
\]
\begin{enumerate}[label=(\alph*)]
\item Show that $\mathbb E[S_N]=\mathbb E[T_N]=0$ and $\mathbb E[S_N^2]=\mathbb E[T_N^2]=\frac N2$.
\item Show that for $N\ge2$,
\[
\mathbb E[S_N^4]=\frac34N^2+\frac98N-3,
\]
so that $\mathbb E[S_N^4]/\mathbb E[S_N^2]^2\to3$, the fourth moment of a standard Gaussian.
\item Show that $\mathbb E[T_N^4]\sim\frac{N^3}3$, so that $\mathbb E[T_N^4]/\mathbb E[T_N^2]^2\to\infty$.
\end{enumerate}
\end{exercise}

\begin{exercise}[$\star\star$]\label{pb:cantor}
\emph{The Fourier transform of the Cantor law.} Let $B_1,B_2,\dots$ be independent fair bits and $X=\sum_{k\ge1}2B_k3^{-k}$, the
Cantor-distributed variable of Tome~3, Exercise~5.7. Let $\varphi(t)=\mathbb E[e^{itX}]$.
\begin{enumerate}[label=(\alph*)]
\item Show that $\varphi(t)=e^{it/2}\prod_{k\ge1}\cos\frac t{3^k}$.
\item Show that $|\varphi(2\pi3^m)|=|\varphi(2\pi)|$ for every $m\ge0$, and that $|\varphi(2\pi)|=0.3714\ldots>0$. Deduce, by a
Riemann--Lebesgue argument, that $X$ has no density.
\item By contrast, show that $\prod_{k\ge1}\cos\frac t{2^k}=\frac{\sin t}t$ (Viète's product), and deduce that $\sum_{k\ge1}\varepsilon_k2^{-k}$, with
independent fair signs $\varepsilon_k$, is uniform on $[-1,1]$.
\end{enumerate}
\end{exercise}

\begin{exercise}[$\star\star\star$]\label{pb:sincsum}
\emph{When a sum equals its integral.} Write $\operatorname{sinc}x=\frac{\sin x}x$, $\operatorname{sinc}0=1$. This is a sequel to
Problem~\ref{pb:sincpower}.
\begin{enumerate}[label=(\alph*)]
\item Show that $\sum_{n\ge1}\frac{\sin n}n=\frac{\pi-1}2$, and deduce $\sum_{n\in\Z}\operatorname{sinc}n=\pi=\int_{-\infty}^\infty\operatorname{sinc}x\dd x$.
\item Let $k\ge2$, $f=\operatorname{sinc}^k$, and let $g_k$ be the density of $U_1+\cdots+U_k$ for independent $U_i$ uniform on $[-1,1]$. Show
that $\hat f(\xi)=\int_{\R}f(x)e^{-2\pi ix\xi}\dd x=2\pi g_k(2\pi\xi)$, and that $g_k$ vanishes outside $[-k,k]$.
\item Deduce that $\sum_{n\in\Z}\operatorname{sinc}^kn=\int_{\R}\operatorname{sinc}^kx\dd x$ for $2\le k\le6$, and that for $k=7$ the sum exceeds the
integral by
\[
\frac{4\pi(7-2\pi)^6}{2^7\cdot6!}=1.849\ldots\cdot10^{-5}.
\]
\end{enumerate}
\end{exercise}

\begin{figure}[ht]
\centering
\fig{ch10-sinc}
\caption{The densities $g_5$, $g_6$, $g_7$ of a sum of $5$, $6$, $7$ independent uniforms on $[-1,1]$. By Problem~\ref{pb:sincsum}, the sum
$\sum_n\operatorname{sinc}^kn$ differs from the integral by $2\pi\sum_{m\ne0}g_k(2\pi m)$. For $k\le6$ the density vanishes beyond $\pm6<2\pi$
(dashed lines); $g_7$ reaches past $2\pi$ by a sliver, $g_7(2\pi)=1.47\cdot10^{-6}$, and the identity fails.}
\label{fig:sinc}
\end{figure}

\section{The Boolean cube}

On the cube $\{-1,1\}^n$, with $x$ uniform, the \emph{characters} are the monomials $x^S=\prod_{i\in S}x_i$ for $S\subseteq\{1,\dots,n\}$. They are
orthonormal: $\mathbb E[x^Sx^T]=\mathbb E[x^{S\triangle T}]$ is $1$ if $S=T$ and $0$ otherwise, because $x_i^2=1$ and $\mathbb E[x_i]=0$. There are $2^n$ of them
in a space of dimension $2^n$, so every $f:\{-1,1\}^n\to\R$ has a unique expansion
\[
f=\sum_S\hat f(S)\,x^S,\qquad\hat f(S)=\mathbb E[f(x)x^S],\qquad\mathbb E[f^2]=\sum_S\hat f(S)^2,
\]
a multilinear polynomial; its \emph{degree} is the largest $|S|$ with $\hat f(S)\ne0$. When $f$ takes values in $\{-1,1\}$ we call it
\emph{Boolean}: think of a vote between two candidates, $x_i$ being voter $i$'s choice and $f(x)$ the outcome. Then
$\sum_S\hat f(S)^2=1$, and the numbers $\hat f(S)^2$ form a probability distribution on the subsets, the \emph{spectral distribution}
of $f$; we write $W^k(f)=\sum_{|S|=k}\hat f(S)^2$. We write $x^{\oplus i}$ for $x$ with its $i$-th coordinate flipped, and $\operatorname{Maj}_n$ (for $n$
odd) for the majority function $\operatorname{sgn}(x_1+\cdots+x_n)$.

\begin{exercise}[$\star\star$]\label{pb:influence}
\emph{Influences.} For a Boolean $f$, the \emph{influence} of voter $i$ is $\operatorname{Inf}_i(f)=P\bigl(f(x)\ne f(x^{\oplus i})\bigr)$, the probability that $i$ is
pivotal, and the \emph{total influence} is $I(f)=\sum_i\operatorname{Inf}_i(f)$.
\begin{enumerate}[label=(\alph*)]
\item Show that $\operatorname{Inf}_i(f)=\sum_{S\ni i}\hat f(S)^2$ and $I(f)=\sum_S|S|\hat f(S)^2$. Deduce the Poincaré inequality $\operatorname{Var}f\le I(f)$, with
equality if and only if $f$ has degree at most $1$.
\item Call $f$ \emph{monotone} if changing a vote from $-1$ to $1$ never changes the outcome from $1$ to $-1$. Show that then
$\operatorname{Inf}_i(f)=\hat f(\{i\})$, and deduce that $I(f)\le\sqrt n$ for every monotone $f$.
\item Show that for $n$ odd, $\operatorname{Inf}_i(\operatorname{Maj}_n)=\binom{n-1}{(n-1)/2}2^{1-n}$, that $I(\operatorname{Maj}_n)\sim\sqrt{\frac{2n}\pi}$, and that $W^1(\operatorname{Maj}_n)\to\frac2\pi$.
\end{enumerate}
\end{exercise}

\begin{exercise}[$\star\star\star$]\label{pb:granularity}
\emph{Low degree forces few variables.} Let $f:\{-1,1\}^n\to\{-1,1\}$ have degree $d\ge1$.
\begin{enumerate}[label=(\alph*)]
\item Show that every $\hat f(S)$ is an integer multiple of $2^{1-d}$.
\item Deduce that at most $4^{d-1}$ coefficients $\hat f(S)$ are non-zero.
\item Show that a non-zero multilinear polynomial of degree at most $k$ does not vanish on at least a fraction $2^{-k}$ of the
cube.
\item Call $i$ \emph{relevant} if $\operatorname{Inf}_i(f)>0$. Show that $\operatorname{Inf}_i(f)\ge2^{1-d}$ for every relevant $i$, and deduce that $f$ has at most
$d\,2^{d-1}$ relevant variables.
\end{enumerate}
\end{exercise}

\begin{exercise}[$\star\star$]\label{pb:bent}
\emph{Bent functions.} Identify $\{-1,1\}^n$ with $\mathbb F_2^n$, so that the characters are $x\mapsto(-1)^{a\cdot x}$, $a\in\mathbb F_2^n$, and write
$\hat f(a)=2^{-n}\sum_x f(x)(-1)^{a\cdot x}$ for $f:\mathbb F_2^n\to\{-1,1\}$.
\begin{enumerate}[label=(\alph*)]
\item Show that $\max_a|\hat f(a)|\ge2^{-n/2}$. Call $f$ \emph{bent} if $|\hat f(a)|=2^{-n/2}$ for every $a$; show that bent functions can
exist only for $n$ even.
\item Show that $f(x,y)=(-1)^{x\cdot y}$ on $\mathbb F_2^m\times\mathbb F_2^m$ is bent.
\item The \emph{nonlinearity} of $f$ is its Hamming distance to the nearest function $\pm(-1)^{a\cdot x}$. Show that it equals
$2^{n-1}\bigl(1-\max_a|\hat f(a)|\bigr)$, hence is at most $2^{n-1}-2^{n/2-1}$, with equality exactly for bent functions.
\item Show that a bent function takes the value $-1$ exactly $2^{n-1}\pm2^{n/2-1}$ times.
\end{enumerate}
\end{exercise}

\begin{exercise}[$\star\star\star$]\label{pb:majoritystable}
\emph{Noise stability and Sheppard's formula.} For $-1<\rho<1$, a pair $(x,y)$ in $\{-1,1\}^n$ is \emph{$\rho$-correlated} if the pairs
$(x_i,y_i)$ are independent, with $x_i$, $y_i$ uniform and $\mathbb E[x_iy_i]=\rho$. The \emph{noise stability} of $f$ is
$\operatorname{Stab}_\rho(f)=\mathbb E[f(x)f(y)]$.
\begin{enumerate}[label=(\alph*)]
\item Show that $\operatorname{Stab}_\rho(f)=\sum_S\rho^{|S|}\hat f(S)^2$.
\item (Sheppard, 1899.) Let $(X,Y)$ be a centred Gaussian pair with unit variances and correlation $\rho$. Show that
$P(X>0,\,Y>0)=\frac14+\frac{\arcsin\rho}{2\pi}$, so that $\mathbb E[\operatorname{sgn}X\operatorname{sgn}Y]=\frac2\pi\arcsin\rho$.
\item Deduce that $\operatorname{Stab}_\rho(\operatorname{Maj}_n)\to\frac2\pi\arcsin\rho$ as $n\to\infty$ through odd values. (You may use the central limit theorem for
random vectors in $\R^2$, stated in Tome~3 after Theorem~6.15.)
\item Deduce that for every $k$, $W^k(\operatorname{Maj}_n)$ converges, to $0$ if $k$ is even and to $\frac2\pi\binom{2j}j\frac1{4^j(2j+1)}$ if $k=2j+1$; check
$W^1\to\frac2\pi$ against Problem~\ref{pb:influence}.
\end{enumerate}
\end{exercise}

\begin{exercise}[$\star\star\star$]\label{pb:arrow}
\emph{Condorcet, Arrow, and Kalai's formula.} Three candidates $a$, $b$, $c$; each of $n$ voters ranks them, choosing one of the six
rankings uniformly at random, independently of the others. A Boolean $f$ decides each of the three duels: let $x_i=1$ if voter $i$
prefers $a$ to $b$ and $x_i=-1$ otherwise, and similarly $y_i$ for $b$ against $c$ and $z_i$ for $c$ against $a$; society prefers $a$ to
$b$ if $f(x)=1$, and so on. The outcome is \emph{rational} if it is not a cycle, that is, unless $f(x)=f(y)=f(z)$.
\begin{enumerate}[label=(\alph*)]
\item Show that each $(x_i,y_i,z_i)$ is uniform on the six vectors of $\{-1,1\}^3$ other than $\pm(1,1,1)$, and deduce that each of the
pairs $(x,y)$, $(y,z)$, $(z,x)$ is $(-\frac13)$-correlated.
\item (Kalai, 2002.) Show that
\[
P(\text{rational})=\frac34-\frac34\operatorname{Stab}_{-1/3}(f).
\]
\item Show that for majority with $n=3$ this probability is $\frac{17}{18}$, and that for $\operatorname{Maj}_n$ it tends to
$\frac34+\frac3{2\pi}\arcsin\frac13=0.9123\ldots$ (Guilbaud, 1952).
\item (Arrow's theorem for three candidates.) Suppose that the outcome is rational for every profile, and that $f$ is
\emph{unanimous}: $f(1,\dots,1)=1$. Show that $f$ is a dictatorship, $f(x)=x_i$ for some $i$.
\end{enumerate}
\end{exercise}

\begin{figure}[ht]
\centering
\fig{ch10-condorcet}
\caption{The probability that majority rule gives a rational outcome with three candidates and $n$ random voters, computed
exactly from Kalai's formula (Problem~\ref{pb:arrow}), for $n=1,3,\dots,51$. It starts at $1$ (one voter is a dictator), drops to
$\frac{17}{18}$ for three voters, and decreases to Guilbaud's limit $0.9123$: about one election in eleven has no Condorcet winner.}
\label{fig:condorcet}
\end{figure}

\begin{exercise}[$\star\star\star$]\label{pb:bonami}
\emph{The Bonami lemma.} Let $f:\{-1,1\}^n\to\R$ be a multilinear polynomial of degree at most $d$.
\begin{enumerate}[label=(\alph*)]
\item Show that $\mathbb E[f^4]\le9^d\,\mathbb E[f^2]^2$.
\item Deduce that if $f\ne0$, then $P\bigl(|f(x)|\ge\frac12\|f\|_2\bigr)\ge\frac9{16}\,9^{-d}$, where $\|f\|_2=\mathbb E[f^2]^{1/2}$.
\item Show that $9$ cannot be replaced by any constant $C<3$: for every $d$ and every $C<3$ there is $f$ of degree $d$ with
$\mathbb E[f^4]>C^d\,\mathbb E[f^2]^2$.
\end{enumerate}
\end{exercise}

\begin{exercise}[Journey]\label{pb:hypercubemix}
\emph{Cutoff on the hypercube.} A point $X_t\in\{-1,1\}^n$ starts at $X_0=(1,\dots,1)$; at each step a coordinate is chosen uniformly
and replaced by a fresh fair sign. Let $d_n(t)$ be the total variation distance between the law of $X_t$ and the uniform law
(Tome~3, Chapter~8).
\begin{enumerate}[label=(\alph*)]
\item Show that $\mathbb E[X_t^S]=\bigl(1-\frac{|S|}n\bigr)^t$ for every $S$.
\item Show that
\[
d_n(t)\le\frac12\Bigl(\sum_{k=1}^n\binom nk\Bigl(1-\frac kn\Bigr)^{2t}\Bigr)^{1/2}\le\frac12\sqrt{\exp\bigl(ne^{-2t/n}\bigr)-1}.
\]
\item Let $W_t$ be the sum of the coordinates of $X_t$. Compute $\mathbb E[W_t]$, show that $\operatorname{Var}W_t\le n$, and use $W_t$ as a test statistic to
bound $d_n(t)$ from below.
\item Deduce the \emph{cutoff}: for every $c>0$,
\[
d_n\bigl(\bigl\lceil\tfrac12n\log n+cn\bigr\rceil\bigr)\le\frac12\sqrt{e^{e^{-2c}}-1},\qquad\liminf_{n\to\infty}d_n\bigl(\bigl\lfloor\tfrac12n\log n-cn\bigr\rfloor\bigr)\ge1-8e^{-2c}.
\]
\end{enumerate}
\end{exercise}

\begin{figure}[ht]
\centering
\fig{ch10-cutoff}
\caption{The exact distance $d_n(t)$ to uniform for the walk of Problem~\ref{pb:hypercubemix}, against $\frac t{n\log n}$, for $n=25$,
$100$ and $400$. As $n$ grows, the curve sharpens into a cliff at $\frac12$: the walk is far from uniform until about $\frac12n\log n$
steps, and close to it soon after.}
\label{fig:cutoff}
\end{figure}

\section{Classical Fourier series}

We come back to the circle of Tome~4: $c_k(f)=\frac1{2\pi}\int_{-\pi}^\pi f(t)e^{-ikt}\dd t$, Parseval's identity (Theorem~6.11), and
trigonometric polynomials $T(\theta)=\sum_{|k|\le n}c_ke^{ik\theta}$, of degree at most $n$. Such a $T$ is real-valued exactly when $c_{-k}=\overline{c_k}$, and
then $z^nT$ restricted to $z=e^{i\theta}$ is a polynomial of degree at most $2n$ in $z$, so a non-zero $T$ has at most $2n$ zeros in a period.

\begin{exercise}[$\star\star\star$]\label{pb:fejer}
\emph{Non-negative trigonometric polynomials.} Let $T(\theta)=\sum_{|k|\le n}c_ke^{ik\theta}$ be real and $T\ge0$.
\begin{enumerate}[label=(\alph*)]
\item (Fejér--Riesz, 1916.) Show that $T(\theta)=|P(e^{i\theta})|^2$ for a polynomial $P\in\C[z]$ of degree at most $n$.
\item (Fejér.) Suppose $c_0=1$. Show that $T(0)\le n+1$, with equality only for the Fejér kernel $\frac1{n+1}\bigl|\sum_{k=0}^ne^{ik\theta}\bigr|^2$.
\item (Egerváry--Szász.) Suppose $T(\theta)=1+\sum_{k=1}^na_k\cos k\theta$ with real $a_k$. Show that $a_1\le2\cos\frac\pi{n+2}$, and that equality is
attained.
\end{enumerate}
\end{exercise}

\begin{exercise}[$\star\star\star$]\label{pb:bernstein}
\emph{Bernstein's inequality.} Let $T$ be a real trigonometric polynomial of degree at most $n$ and $\|T\|=\max|T|$.
\begin{enumerate}[label=(\alph*)]
\item Show that if $\|T\|\le1$, then $T'(\theta)^2+n^2T(\theta)^2\le n^2$ for all $\theta$. Deduce $\|T'\|\le n\|T\|$, and that the same holds for
complex-valued $T$.
\item Deduce that a real polynomial $P$ of degree $n$ with $|P|\le1$ on $[-1,1]$ satisfies $|P'(x)|\le\frac n{\sqrt{1-x^2}}$ for $|x|<1$.
\item Show that both inequalities are sharp.
\end{enumerate}
\end{exercise}

\begin{exercise}[$\star\star\star$]\label{pb:bernsteinabs}
\emph{Absolutely convergent Fourier series.} Let $f$ be $2\pi$-periodic with $|f(x)-f(y)|\le C|x-y|^\alpha$ for all $x,y$, where
$0<\alpha\le1$.
\begin{enumerate}[label=(\alph*)]
\item Show that for every $h$, $\frac1{2\pi}\int_{-\pi}^\pi|f(x+h)-f(x-h)|^2\dd x=\sum_k4\sin^2(kh)\,|c_k(f)|^2$.
\item (S.~Bernstein, 1914.) Suppose $\alpha>\frac12$. Show that $\sum_k|c_k(f)|<\infty$, and that the Fourier series of $f$ converges
uniformly to $f$.
\item Let $P_j$ be the Rudin--Shapiro polynomials of Problem~\ref{pb:rudin}, and $f(x)=\sum_{j\ge0}2^{-j}e^{i2^jx}P_j(e^{ix})$. Show that $f$
satisfies the hypothesis with $\alpha=\frac12$, but that $\sum_k|c_k(f)|=\infty$. So $\frac12$ is the threshold.
\end{enumerate}
\end{exercise}

\begin{exercise}[$\star\star$]\label{pb:fantaussky}
\emph{Polygons by the discrete Fourier transform.} Let $n\ge3$ and $\omega=e(\frac1n)$; indices are taken modulo $n$.
\begin{enumerate}[label=(\alph*)]
\item (Fan, Taussky and Todd, 1955.) Let $x_0,\dots,x_{n-1}\in\C$ with $\sum_kx_k=0$. Show that
\[
\sum_k|x_{k+1}-x_k|^2\ge4\sin^2\frac\pi n\sum_k|x_k|^2,
\]
with equality if and only if $x_k=a\omega^k+b\omega^{-k}$ for some $a,b\in\C$.
\item Let $z_0,\dots,z_{n-1}$ be the vertices of a closed polygon in $\C$, with signed area $A=\frac12\operatorname{Im}\sum_k\overline{z_k}z_{k+1}$. Show that
\[
\sum_k|z_{k+1}-z_k|^2\ge4\tan\frac\pi n\cdot A,
\]
with equality if and only if the polygon is a positively oriented regular $n$-gon (or a single point).
\item Deduce that among equilateral $n$-gons of given perimeter, the regular one has the largest area.
\end{enumerate}
\end{exercise}

\begin{exercise}[$\star\star\star$]\label{pb:siegel}
\emph{Minkowski's theorem by Fourier analysis (Siegel, 1935).} Let $K\subset\R^n$ be a bounded convex set with $K=-K$, and suppose
that $K$ contains no point of $\Z^n$ other than $0$.
\begin{enumerate}[label=(\alph*)]
\item Show that the translates $\frac12K+m$, $m\in\Z^n$, are pairwise disjoint.
\item Let $\varphi(x)=\sum_{m\in\Z^n}1_{\frac12K}(x+m)$, a $\Z^n$-periodic function. Show that $\varphi^2=\varphi$, and that its Fourier coefficients
on the torus $[0,1[^n$ are $\hat\varphi(0)=2^{-n}\operatorname{vol}K$ and $\hat\varphi(k)=\int_{\frac12K}e(-k\cdot x)\dd x$.
\item Using Bessel's inequality on the torus, deduce that
\[
2^{-n}\operatorname{vol}K\bigl(1-2^{-n}\operatorname{vol}K\bigr)\ge\sum_{k\ne0}\Bigl|\int_{\frac12K}e(-k\cdot x)\dd x\Bigr|^2,
\]
and in particular Minkowski's theorem: a symmetric convex body of volume greater than $2^n$ contains a non-zero integer point.
\item Check that for the open cube $K=\left]-1,1\right[^n$ both sides vanish.
\end{enumerate}
\end{exercise}

\section{Hints}

\begin{hint}{pb:bilinear}
(a) Cauchy--Schwarz over $a$, then let $a$ run over all of $\mathbb F_p$ and expand the square. (b) Write $N(x)$ with the orthogonality
relation, isolate $r=0$, apply Cauchy--Schwarz over $v\in V$ and extend the sum to all $v\in\mathbb F_p^2$. In the expanded square, the
pairs $(r,r')$ with $r\ne r'$ contribute sums $\sum_{r'}e(\frac{-r'(t-1)x}p)$ with $t=\frac r{r'}\ne1$. (c) $U=A\times C$, $V=B\times D$. (d) Look at the
real parts.
\end{hint}

\begin{hint}{pb:kloosterman}
(a) Substitute $x\mapsto-x$, and $x\mapsto bx$. (b) Sum over $a$ first. (c) Do not attack $K(a)^4$ directly: compute $\sum_{a,b}|K(a,b)|^4$ in two
ways, once with (a) and once by summing over $a$ and $b$ first, which counts solutions of $x_1+x_2=x_3+x_4$, $\bar x_1+\bar x_2=\bar x_3+\bar x_4$.
(d) One term of a sum of non-negative terms is at most the sum.
\end{hint}

\begin{hint}{pb:apcount}
(a) A triple $(x,y,z)$ is a progression iff $x+z=2y$. (b) Isolate $r=0$ and use Parseval for the rest. (c) Write $1_Q(x)=\frac12(1+\chi(x))$
for $x\ne0$, with $\chi$ the Legendre symbol, and use the Gauss sum. (d) Progressions with $d=0$ number $|A|$.
\end{hint}

\begin{hint}{pb:dftspectrum}
(a) $(F^2x)_j=\frac1n\sum_kx_k\sum_le(-\frac{l(j+k)}n)$. A matrix killed by a polynomial with simple roots is diagonalisable. (b) $\operatorname{tr}F$ is a
quadratic Gauss sum, whose value Gauss determined (Tome~4, Theorem~8.14). (c) Four linear equations in four unknowns.
\end{hint}

\begin{hint}{pb:newman}
(a) $\frac13\sum_{t=0}^2\omega^{t(s(n)-j)}$ is the indicator of $s(n)\equiv j$; the sum of $\omega^{ts(n)}$ over $n<2^k$ factorises over the $k$ bits.
(b) Now filter $3\mid n$ with $\frac13\sum_t\omega^{tn}$, and note that $2^i\bmod3$ alternates between $1$ and $2$.
\end{hint}

\begin{hint}{pb:erdosturan}
(a) The Fourier coefficients of a convolution multiply. (b) Choose $J$ so that $t\ge1_I$ for an upper bound, and $t\le1_I$ for a lower
bound; expand $t$ in its (absolutely convergent) Fourier series and sum over the points. (c) Gauss sums have modulus $\sqrt p$;
take $\delta=p^{-1/2}$.
\end{hint}

\begin{hint}{pb:circlewalk}
$\mathbb E[e(hX_n)]=\cos^n(2\pi h\alpha)$. When is $|\cos(2\pi h\alpha)|=1$?
\end{hint}

\begin{hint}{pb:kac}
(a) Products of cosines are sums of cosines. (b) $\mathbb E[S_N^4]$ is $\frac1{16}$ of the number of $(k_1,\dots,k_4,\varepsilon_1,\dots,\varepsilon_4)$ with
$\sum\varepsilon_i2^{k_i}=0$. By uniqueness of binary expansions, such a relation either cancels in pairs or reads $2^{m+2}=2^{m+1}+2^m+2^m$.
(c) The same count for $\sum\varepsilon_ik_i=0$ is a volume.
\end{hint}

\begin{hint}{pb:cantor}
(a) Independence turns the characteristic function of a sum into a product. (b) Multiplying $t$ by $3$ shifts the product by one
factor, and the extra factor is $\cos(2\pi\cdot\text{integer})$. For Riemann--Lebesgue, approximate a density by step functions.
(c) $\sin t=2^m\sin\frac t{2^m}\prod_{k=1}^m\cos\frac t{2^k}$.
\end{hint}

\begin{hint}{pb:sincsum}
(a) The sawtooth $\frac{\pi-x}2$ on $]0,2\pi[$. (b) $\operatorname{sinc}x=\frac12\int_{-1}^1e^{ixu}\dd u$ is the characteristic function of a uniform variable.
(c) Poisson summation (Tome~4, Theorem~8.9).
\end{hint}

\begin{hint}{pb:influence}
(a) $f(x)\ne f(x^{\oplus i})$ iff the derivative $D_if(x)=\frac12\bigl(f(x^{i\to1})-f(x^{i\to-1})\bigr)$ is $\pm1$. (b) For monotone $f$, $D_if\in\{0,1\}$.
Cauchy--Schwarz. (c) Voter $i$ is pivotal iff the others split evenly. Stirling.
\end{hint}

\begin{hint}{pb:granularity}
(a) Write $g=\frac{1-f}2$ as a polynomial in the variables $y_i=\frac{1-x_i}2\in\{0,1\}$: its coefficients are integers. (c) Induction on $n$:
write $p=x_nq+r$. (d) Apply (c) to the derivative $D_if$, and use $I(f)=\sum|S|\hat f(S)^2\le d$.
\end{hint}

\begin{hint}{pb:bent}
(a) Parseval, and integrality of $2^n\hat f(a)$. (b) Sum over $y$ first. (c) The distance from $f$ to $g=\pm(-1)^{a\cdot x}$ is
$2^{n-1}(1-\mathbb E[fg])$. (d) Look at $\hat f(0)$.
\end{hint}

\begin{hint}{pb:majoritystable}
(a) $\mathbb E[x^Sy^T]$. (b) Write $Y=\rho X+\sqrt{1-\rho^2}Z$ and use the rotation invariance of $(X,Z)$: the event is a wedge. (c) $(\frac{\sum x_i}{\sqrt n},\frac{\sum y_i}{\sqrt n})$.
(d) Power series with non-negative coefficients of sum $1$ that converge pointwise on $[0,1[$; the Taylor series of $\arcsin$.
\end{hint}

\begin{hint}{pb:arrow}
(a) List the six rankings. (b) For $u,v,w\in\{-1,1\}$, $1-1_{u=v=w}=\frac34-\frac14(uv+vw+wu)$. (c) For odd $f$, $\operatorname{Stab}_{-\rho}(f)=-\operatorname{Stab}_\rho(f)$.
(d) $(-\frac13)^k\ge-\frac13$, with equality only for $k=1$. A Boolean function of degree $1$ is $\pm x_i$.
\end{hint}

\begin{hint}{pb:bonami}
(a) Induction on $n$: $f=x_ng+h$ with $g$, $h$ not involving $x_n$, $\deg g\le d-1$. Expand the fourth power and use Cauchy--Schwarz on the
middle term. (b) Paley--Zygmund: Cauchy--Schwarz on $\mathbb E[f^21_{f^2\ge s/4}]$. (c) Products of normalised sums over disjoint blocks.
\end{hint}

\begin{hint}{pb:hypercubemix}
(a) One step multiplies $\mathbb E[X^S]$ by $1-\frac{|S|}n$. (b) Write the law of $X_t$ as $2^{-n}h_t$ and bound $\mathbb E|h_t-1|$ by $\mathbb E[(h_t-1)^2]^{1/2}$; use
$\binom nk\le\frac{n^k}{k!}$ and $1-u\le e^{-u}$. (c) $W_t=\sum_ix_i$; compute $\mathbb E[W_t^2]$ with (a). Under the uniform law, $W$ has mean $0$ and variance
$n$; use Chebyshev twice.
\end{hint}

\begin{hint}{pb:fejer}
(a) Look at the roots of the polynomial $z^nT(z)$: they come in pairs $\alpha$, $1/\bar\alpha$, and those on the unit circle have even
multiplicity. (b) Cauchy--Schwarz on $P(1)$. (c) $a_1=2\operatorname{Re}\sum p_kp_{k+1}$ is a quadratic form; the adjacency matrix of a path.
\end{hint}

\begin{hint}{pb:bernstein}
(a) If the inequality fails at $\theta_0$, choose $\beta$ with $T(\theta_0)=\cos(n(\theta_0-\beta))$ and count the zeros of $\cos(n(\theta-\beta))-T(\theta)$. (b) $T(\theta)=P(\cos\theta)$.
(c) $\sin n\theta$ and the Chebyshev polynomials.
\end{hint}

\begin{hint}{pb:bernsteinabs}
(a) Parseval for $f(x+h)-f(x-h)$. (b) Dyadic blocks $2^{j-1}<|k|\le2^j$ with $h=\frac\pi{3\cdot2^j}$, then Cauchy--Schwarz in each block. (c) Split
$f$ at the scale of $h$: Bernstein's inequality for the low blocks, the sup norm for the high ones.
\end{hint}

\begin{hint}{pb:fantaussky}
(a) Expand $x_k=\sum_mc_m\omega^{mk}$ and use Parseval for the DFT. (b) Do the same for the area. (c) For an equilateral polygon,
$L^2=n\sum|z_{k+1}-z_k|^2$.
\end{hint}

\begin{hint}{pb:siegel}
(a) If $x+m=y+m'$ with $x,y\in\frac12K$, then $m'-m\in K$. (b) Unfold the integral over the torus into an integral over $\R^n$. (c) Bessel:
$\int|\varphi|^2\ge\sum_{k\in F}|\hat\varphi(k)|^2$ for every finite $F$.
\end{hint}

\section{Solutions}

\begin{solution}{pb:bilinear}
\idea A bilinear sum is small because one of the two variables can be freed: Cauchy--Schwarz lets it run over all of $\mathbb F_p$,
and then orthogonality kills every off-diagonal term. In two dimensions the same trick counts points on lines.

\writeup (a) Let $\sigma(a)=\sum_{b\in B}e(\frac{rab}p)$. By Cauchy--Schwarz and then by enlarging the range of $a$,
\[
\Bigl|\sum_{a\in A}\sigma(a)\Bigr|^2\le|A|\sum_{a\in A}|\sigma(a)|^2\le|A|\sum_{a\in\mathbb F_p}|\sigma(a)|^2=|A|\sum_{b,b'\in B}\sum_{a\in\mathbb F_p}e\Bigl(\frac{ra(b-b')}p\Bigr).
\]
As $r\ne0$, the inner sum is $p$ if $b=b'$ and $0$ otherwise, so the right side is $p|A||B|$.

(b) By orthogonality,
\[
N(x)=\frac1p\sum_{r\in\mathbb F_p}\sum_{u\in U,v\in V}e\Bigl(\frac{r(u\cdot v-x)}p\Bigr).
\]
The term $r=0$ is $\frac{|U||V|}p$. With $\hat U(w)=\sum_{u\in U}e(\frac{u\cdot w}p)$ for $w\in\mathbb F_p^2$, the rest is
\[
E=\frac1p\sum_{v\in V}\Sigma(v),\qquad\Sigma(v)=\sum_{r\ne0}e\Bigl(-\frac{rx}p\Bigr)\hat U(rv).
\]
By Cauchy--Schwarz, $|E|^2\le\frac{|V|}{p^2}\sum_{v\in V}|\Sigma(v)|^2\le\frac{|V|}{p^2}\sum_{v\in\mathbb F_p^2}|\Sigma(v)|^2$. Expanding the square and summing over $v$ first,
\[
\sum_{v\in\mathbb F_p^2}|\Sigma(v)|^2=\sum_{r,r'\ne0}e\Bigl(-\frac{(r-r')x}p\Bigr)\sum_{u,u'\in U}\sum_ve\Bigl(\frac{v\cdot(ru-r'u')}p\Bigr)=p^2\sum_{r,r'\ne0}e\Bigl(-\frac{(r-r')x}p\Bigr)\nu(r,r'),
\]
where $\nu(r,r')$ is the number of $(u,u')\in U^2$ with $ru=r'u'$. The diagonal $r=r'$ gives $(p-1)|U|$. Off the diagonal, put $t=\frac r{r'}$,
so $t\ne0,1$ and $\nu(r,r')=\#\{u\in U:tu\in U\}=:\mu(t)$ depends only on $t$; for fixed $t$ the sum over $r'\ne0$ of $e(-\frac{r'(t-1)x}p)$
is $-1$, since $(t-1)x\ne0$. So the off-diagonal part is $-\sum_{t\ne0,1}\mu(t)\le0$, and
\[
|E|^2\le\frac{|V|}{p^2}\cdot p^2(p-1)|U|<p|U||V|.
\]

(c) Take $U=A\times C$ and $V=B\times D$: then $N(x)$ counts the $(a,c,b,d)$ with $ab+cd=x$, and by (b)
$N(x)\ge\frac{|A||B||C||D|}p-\sqrt{p|A||B||C||D|}$, which is positive exactly when $|A||B||C||D|>p^3$. With $A=B=C=D$ the condition is $|A|^4>p^3$.

(d) For $1\le a,b\le m$, $0<\frac{ab}p\le\frac{m^2}p\le\frac19$, so $\cos\frac{2\pi ab}p\ge\cos\frac{2\pi}9=0.766\ldots>\frac12$ and the real part of the sum is at
least $\frac{m^2}2$. With $m=\lfloor\frac{\sqrt p}3\rfloor$ we have $\sqrt{p|A||B|}=\sqrt p\,m\le(3m+3)m$, so the sum is at least
$\frac{m}{6(m+1)}\sqrt{p|A||B|}$, which is at least $\frac1{12}\sqrt{p|A||B|}$ once $p\ge9$.

\tools{Tome~4, Theorem~7.3 (orthogonality and Parseval for the DFT); Tome~0, Chapter~1 (Cauchy--Schwarz).}
\leads Bound (a) is the starting point of \emph{sum--product} theory over $\mathbb F_p$: Bourgain, Katz and Tao (2004) proved that a set
$A\subseteq\mathbb F_p$ with $p^\varepsilon<|A|<p^{1-\varepsilon}$ has $|A+A|+|A\cdot A|\ge|A|^{1+\delta}$, the finite-field analogue of Problem~\ref{pb:sumproduct}. Part~(b) is an
incidence bound between points and lines of the plane $\mathbb F_p^2$, a finite-field counterpart of Problem~\ref{pb:incidences}.
\end{solution}

\begin{solution}{pb:kloosterman}
\idea Second and fourth moments of a family of exponential sums are counts of solutions of equations, which are easy. Since a
single term is at most the whole fourth moment, a bound for the moment is a bound for each sum.

\writeup (a) The substitution $x\mapsto-x$ permutes $\mathbb F_p^\times$ and changes $\bar x$ into $-\bar x$, so it turns $K(a)$ into $\overline{K(a)}$: hence $K(a)$
is real. $K(0)=\sum_{x\ne0}e(\frac{\bar x}p)=\sum_{y\ne0}e(\frac yp)=-1$. For $b\ne0$, substitute $x=by$: then $b\bar x=\bar y$ and $ax=aby$, so $K(a,b)=K(ab,1)$.

(b) Expanding, and summing over $a$ first,
\[
\sum_aK(a)\overline{K(a)}=\sum_{x,y\ne0}e\Bigl(\frac{\bar x-\bar y}p\Bigr)\sum_ae\Bigl(\frac{a(x-y)}p\Bigr)=p\sum_{x\ne0}1=p(p-1).
\]

(c) Let $M=\sum_aK(a)^4$. Compute $\Sigma=\sum_{a,b\in\mathbb F_p}|K(a,b)|^4$ in two ways. First, for $b\ne0$, (a) and the bijection $a\mapsto ab$ give
$\sum_a|K(a,b)|^4=M$; for $b=0$, $K(a,0)=\sum_{x\ne0}e(\frac{ax}p)$ is $p-1$ for $a=0$ and $-1$ otherwise. So
$\Sigma=(p-1)M+(p-1)^4+(p-1)$. Second, $|K(a,b)|^4=K(a,b)^2\,\overline{K(a,b)}^2$ expands into
$\sum_{x_1,\dots,x_4}e\bigl(\frac{a(x_1+x_2-x_3-x_4)+b(\bar x_1+\bar x_2-\bar x_3-\bar x_4)}p\bigr)$, and summing over $a$ and $b$ gives $\Sigma=p^2N$, where $N$ counts the
$(x_1,\dots,x_4)\in(\mathbb F_p^\times)^4$ with
\[
x_1+x_2=x_3+x_4,\qquad\bar x_1+\bar x_2=\bar x_3+\bar x_4.
\]
Let $s=x_1+x_2$. If $s\ne0$, then $\bar x_1+\bar x_2=\frac s{x_1x_2}$, so the equations say $x_3+x_4=s$ and $x_3x_4=x_1x_2$: $x_3$ and $x_4$ are the
roots of $X^2-sX+x_1x_2$, that is $\{x_3,x_4\}=\{x_1,x_2\}$, which gives $2$ solutions if $x_1\ne x_2$ and $1$ if $x_1=x_2$. There are $(p-1)(p-2)$
pairs $(x_1,x_2)$ with $s\ne0$, of which $p-1$ have $x_1=x_2$. If $s=0$, then $\bar x_1+\bar x_2=0$ too, and $(x_3,x_4)$ is any pair with
$x_4=-x_3$: $(p-1)^2$ solutions in all. So
\[
N=2(p-1)(p-2)-(p-1)+(p-1)^2=3(p-1)(p-2).
\]
Comparing, $(p-1)M=3p^2(p-1)(p-2)-(p-1)^4-(p-1)$, and dividing by $p-1$,
\[
M=3p^2(p-2)-(p-1)^3-1=2p^3-3p^2-3p.
\]

(d) Every term of $M$ is non-negative, so $K(a)^4\le M<2p^3$ and $|K(a)|<2^{1/4}p^{3/4}<1.19\,p^{3/4}$. The trivial bound is $p-1$, and (b) gives
only $|K(a)|<p$, no better. The fourth moment wins as soon as $1.19\,p^{3/4}<p-1$, that is for every $p\ge5$, and by a growing factor.

\tools{Tome~1, Chapter~3 (inverses modulo $p$); Tome~0, Chapter~3 (a quadratic has at most two roots in a field).}
\leads Kloosterman introduced these sums in 1926 to count representations of integers by quaternary quadratic forms, and his
exponent $\frac34$ was enough for that purpose. In 1948 André Weil deduced from his proof of the Riemann hypothesis for curves over
finite fields that $|K(a)|\le2\sqrt p$ for $a\ne0$, the optimal exponent. Writing $K(a)=2\sqrt p\cos\theta_a$, Katz
(1988) showed that the angles $\theta_a$ are distributed, as $p\to\infty$, with density $\frac2\pi\sin^2\theta$ on $[0,\pi]$: the Sato--Tate law,
which can be observed numerically.
\end{solution}

\begin{solution}{pb:apcount}
\idea A progression is a solution of the linear equation $x+z=2y$, so its count is a Fourier sum. The trivial frequency gives the
number of progressions in a random set of the same density, and the others are controlled by the largest non-trivial
coefficient: small coefficients mean random-looking, and random-looking sets have the expected number of progressions.

\writeup (a) The map $(x,d)\mapsto(x,x+d,x+2d)$ is a bijection from $\mathbb F_p^2$ onto the triples $(x,y,z)$ with $x+z=2y$ (given such a triple,
$d=y-x$ and then $z=2y-x=x+2d$). So
\[
T(A)=\sum_{x,y,z\in A}\frac1p\sum_re\Bigl(\frac{r(x+z-2y)}p\Bigr)=\frac1p\sum_r\hat A(r)^2\hat A(-2r).
\]

(b) The term $r=0$ is $\frac{|A|^3}p$. For $r\ne0$, $-2r\ne0$ since $p$ is odd, so the rest is at most
$\frac1p\max_{s\ne0}|\hat A(s)|\sum_r|\hat A(r)|^2=|A|\max_{s\ne0}|\hat A(s)|$ by Parseval.

(c) Let $\chi$ be the Legendre symbol, $\chi(0)=0$. For $x\ne0$, $1_Q(x)=\frac12(1+\chi(x))$, so for $r\ne0$,
\[
\hat Q(r)=\frac12\sum_{x\ne0}e\Bigl(\frac{rx}p\Bigr)+\frac12\sum_x\chi(x)e\Bigl(\frac{rx}p\Bigr)=\frac{-1+\chi(r)g}2,
\]
where $g=\sum_x\chi(x)e(\frac xp)$ is the Gauss sum (substitute $x\mapsto\bar rx$), and $|g|=\sqrt p$. Hence $|\hat Q(r)|\le\frac{\sqrt p+1}2$, and (b) with
$|Q|=\frac{p-1}2$ gives the claim.

(d) By (b), $T(A)\ge\delta^3p^2-\delta p\cdot\frac12\delta^2p=\frac12\delta^3p^2$. The pairs with $d=0$ number $|A|=\delta p$, and $\frac12\delta^3p^2>\delta p$ because
$\delta^2p>2$. So some pair has $d\ne0$.

\tools{Tome~4, Theorem~7.3 and Proposition~7.13 (Parseval, the Gauss sum); Tome~1, Chapter~6 (the Legendre symbol).}
\leads Part~(d) is the first half of Roth's proof (1953) that a subset of $\{1,\dots,N\}$ of density $\delta$ contains a three-term
progression once $N$ is large: either all non-trivial coefficients are small, and (d) finds a progression, or one is large, and
the set is denser on a long subprogression, where one starts again. Meshulam's version over $\mathbb F_3^n$ is in Tome~7. Behrend (1946)
built progression-free sets of size $Ne^{-c\sqrt{\log N}}$, and Kelley and Meka (2023) proved a nearly matching upper bound of the
shape $Ne^{-c(\log N)^\beta}$.
\end{solution}

\begin{solution}{pb:dftspectrum}
\idea $F^4=I$ pins the spectrum to four points, and two traces, one of which is Gauss's famous sum, pin the multiplicities.

\writeup (a) By orthogonality,
\[
(F^2x)_j=\frac1n\sum_le\Bigl(-\frac{jl}n\Bigr)\sum_kx_ke\Bigl(-\frac{lk}n\Bigr)=\sum_kx_k\cdot\frac1n\sum_le\Bigl(-\frac{l(j+k)}n\Bigr)=x_{-j}.
\]
So $F^2=R$, the reversal, and $F^4=R^2=I$. The polynomial $X^4-1$ kills $F$ and has four simple roots, so $F$ is diagonalisable (its
minimal polynomial divides $X^4-1$, hence has simple roots), with eigenvalues among $1,-1,-i,i$.

(b) In an eigenbasis,
\[
\operatorname{tr}F=m_1-m_{-1}-im_{-i}+im_i,\qquad\operatorname{tr}F^2=m_1+m_{-1}-m_{-i}-m_i.
\]
Now $\operatorname{tr}F^2=\operatorname{tr}R$ is the number of $k\in\Z/n\Z$ with $k=-k$: $1$ if $n$ is odd and $2$ if $n$ is even. And
$\operatorname{tr}F=\frac1{\sqrt n}\sum_ke(-\frac{k^2}n)=\frac{\overline{G(n)}}{\sqrt n}$, where $G(n)=\sum_{k=0}^{n-1}e(\frac{k^2}n)$; by Gauss's theorem $\operatorname{tr}F$ is $1-i$, $1$, $0$, $-i$ for
$n\equiv0,1,2,3\pmod4$.

(c) Let $\sigma=m_1+m_{-1}$ and $\tau=m_{-i}+m_i$. Then $\sigma+\tau=n$ and $\sigma-\tau=\operatorname{tr}F^2$; and $m_1-m_{-1}=\operatorname{Re}\operatorname{tr}F$, $m_i-m_{-i}=\operatorname{Im}\operatorname{tr}F$. For
$n=4m$: $\sigma=2m+1$, $\tau=2m-1$, $m_1-m_{-1}=1$, $m_i-m_{-i}=-1$, so $(m_1,m_{-1},m_{-i},m_i)=(m+1,m,m,m-1)$. For $n=4m+1$: $\sigma=2m+1$, $\tau=2m$,
differences $1$ and $0$: $(m+1,m,m,m)$. For $n=4m+2$: $\sigma=2m+2$, $\tau=2m$, differences $0$, $0$: $(m+1,m+1,m,m)$. For $n=4m+3$: $\sigma=2m+2$,
$\tau=2m+1$, differences $0$ and $-1$: $(m+1,m+1,m+1,m)$.

\tools{Tome~4, Theorem~7.3 (DFT inversion) and Theorem~8.14 (Gauss's evaluation); Tome~0, Chapter~7 (diagonalisation).}
\leads The argument can be run backwards: Schur (1921) found the multiplicities independently, from the determinant of the DFT
matrix (a Vandermonde determinant), and deduced the value of $G(n)$, sign included, which had cost Gauss years of effort. On $L^2(\R)$ the
Fourier transform also satisfies $\mathcal F^4=I$; its eigenfunctions are the Hermite functions $H_k(\sqrt{2\pi}\,x)e^{-\pi x^2}$, with eigenvalue $(-i)^k$, and
McClellan and Parks (1972) studied discrete analogues of them for the DFT.
\end{solution}

\begin{solution}{pb:newman}
\idea To pick out a residue class modulo $3$, average over the three characters of $\Z/3\Z$. Digit sums then factorise over the
bits, and everything reduces to powers of $1+\omega$ and $1-\omega$.

\writeup (a) Let $N_j$ be the count. Since $\frac13\sum_{t=0}^2\omega^{tm}$ is $1$ if $3\mid m$ and $0$ otherwise,
\[
N_j=\frac13\sum_{t=0}^2\omega^{-tj}\sum_{n<2^k}\omega^{ts(n)}=\frac13\sum_{t=0}^2\omega^{-tj}(1+\omega^t)^k,
\]
because $n\mapsto$ its $k$ bits is a bijection onto $\{0,1\}^k$ and $\omega^{ts(n)}$ is the product over the bits of $\omega^{t\cdot\text{bit}}$. The term $t=0$ is
$2^k$; $1+\omega=\frac12+\frac{\sqrt3}2i=e^{i\pi/3}$, so the term $t=1$ is $e^{i\pi(k-2j)/3}$, and the term $t=2$ is its conjugate. Hence
$N_j=\frac13\bigl(2^k+2\cos\frac{\pi(k-2j)}3\bigr)$.

(b) Filtering the multiples of $3$ in the same way, for $K\ge1$,
\[
\sum_{\substack{n<2^K\\3\mid n}}(-1)^{s(n)}=\frac13\sum_{t=0}^2\sum_{n<2^K}\omega^{tn}(-1)^{s(n)}=\frac13\sum_{t=0}^2\prod_{i=0}^{K-1}\bigl(1-\omega^{t2^i}\bigr),
\]
since $n=\sum b_i2^i$ gives $\omega^{tn}(-1)^{s(n)}=\prod_i(-\omega^{t2^i})^{b_i}$. The term $t=0$ vanishes. As $2^i\equiv1$ or $2\pmod3$ according to the parity of $i$,
the factors for $t=1$ alternate between $1-\omega$ and $1-\omega^2$, and $(1-\omega)(1-\omega^2)=2-\omega-\omega^2=3$. For $K=2k$ the terms $t=1,2$ are
both $3^k$, and the sum is $\frac23\cdot3^k=2\cdot3^{k-1}$. For $K=2k+1$ they are $3^k(1-\omega)$ and $3^k(1-\omega^2)$, and the sum is $\frac13\cdot3^k\cdot3=3^k$.
With $N=4^k$ the margin is $3^k=N^{\log3/\log4}=N^{0.79\ldots}$.

\tools{Tome~0, Chapter~2 (roots of unity); Tome~4, Chapter~7 (characters of $\Z/n\Z$).}
\leads D.~J.~Newman (1969) proved that the sum $\sum_{n<N,\,3\mid n}(-1)^{s(n)}$ is positive for \emph{every} $N\ge1$, and of exact order
$N^{\log3/\log4}$. J.~Coquet (1983) found the exact formula: the sum is $N^{\log3/\log4}F(\log_4N)$ plus a bounded error, with $F$ continuous,
periodic and nowhere differentiable. Similar biases hold for other moduli and bases, and Mauduit and Rivat (2010) proved that
$s(p)$ is equidistributed modulo $2$ along the primes (Problem~\ref{pb:digitsclt}).
\end{solution}

\begin{solution}{pb:erdosturan}
\idea Indicators of intervals have badly decaying Fourier coefficients, so we sandwich them between continuous trapezoids, whose
coefficients decay like $h^{-2}$. The price is the width $\delta$ of the slopes; the gain is an absolutely convergent expansion that
turns counting points into exponential sums.

\writeup (a) Write $k=\frac1\delta1_{[0,\delta]}$, extended $1$-periodically like $1_J$. Then $t(x)=\frac1\delta\int_0^\delta1_J(x-u)\dd u$ is an average of values in
$\{0,1\}$, so $0\le t\le1$, and by Fubini $\int_0^1t=\int_0^11_J=L$. The Fourier coefficients of a periodic convolution multiply (Fubini
again): $\hat t(h)=\hat1_J(h)\hat k(h)$. For $h\ne0$,
\[
|\hat1_J(h)|=\Bigl|\int_Je(-hx)\dd x\Bigr|=\frac{|\sin\pi hL|}{\pi|h|}\le\frac1{\pi|h|},\qquad|\hat k(h)|=\frac{|\sin\pi h\delta|}{\pi|h|\delta}\le\min\Bigl(1,\frac1{\pi|h|\delta}\Bigr),
\]
and the bound follows.

(b) By (a), $\sum_h|\hat t(h)|<\infty$, so $\sum_h\hat t(h)e(hx)$ converges uniformly to a continuous function with the same coefficients as $t$, which is
$t$ itself (Tome~4, Corollary~6.10(b), after the change of variable $x\mapsto\frac x{2\pi}$). Hence
\[
\frac1N\sum_nt(x_n)=L+\frac1N\sum_{h\ne0}\hat t(h)S(h),
\]
and, as $S(-h)=\overline{S(h)}$, the error is at most
\[
E=\frac2N\sum_{h\ge1}\min\Bigl(\frac1{\pi h},\frac1{\pi^2h^2\delta}\Bigr)|S(h)|.
\]
Let $I$ have ends $a<b$ and length $\ell$,
and let $\nu=\frac1N\#\{n:x_n\in I\}$.

Upper bound. If $\ell\le1-2\delta$, take $J=[a-\delta,b]$, of length $L=\ell+\delta\le1-\delta$. For $x\in I$, $[x-\delta,x]\subseteq J$, so $t(x)=1$: thus $t\ge1_I$
and $\nu\le\frac1N\sum t(x_n)\le\ell+\delta+E$. If $\ell>1-2\delta$, simply $\nu\le1<\ell+2\delta$.

Lower bound. If $\ell\ge\delta$, take $J=[a,b-\delta]$, of length $\ell-\delta$. If $t(x)>0$ then $x-u\in J$ for a set of $u\in[0,\delta]$ of positive
length, which forces $a<x<b$; so $t\le1_I$ and $\nu\ge\ell-\delta-E$. If $\ell<\delta$, simply $\nu\ge0>\ell-\delta$.

In all cases $|\nu-\ell|\le2\delta+E$, which is the claim.

(c) For $p\nmid h$, $S(h)=\sum_{k=0}^{p-1}e(\frac{hk^2}p)=\sum_y(1+\chi(y))e(\frac{hy}p)=\chi(h)g$, where $\chi$ is the Legendre symbol and $g$ the Gauss sum; so
$|S(h)|=\sqrt p$. For $p\mid h$, $S(h)=p$. We use two elementary bounds, for $X\ge1$: $\sum_{h\le X}\frac1h\le1+\log X$, and
$\sum_{h>X}\frac1{h^2}\le\sum_{h>\lfloor X\rfloor}\frac1{h(h-1)}=\frac1{\lfloor X\rfloor}\le\frac2X$. Take $\delta=p^{-1/2}$. The terms with $p\nmid h$ contribute at most
\[
\frac{2\sqrt p}p\Bigl(\sum_{h\le\sqrt p}\frac1{\pi h}+\sum_{h>\sqrt p}\frac1{\pi^2h^2\delta}\Bigr)\le\frac2{\sqrt p}\Bigl(\frac{1+\frac12\log p}\pi+\frac2{\pi^2}\Bigr),
\]
and those with $h=mp$ at most $\frac2p\sum_{m\ge1}\frac p{\pi^2m^2p^2\delta}=\frac1{3p^{3/2}}$. Altogether
\[
D\le\frac{\log p+2+2\pi+\frac4\pi}{\pi\sqrt p}+\frac1{3p^{3/2}}\le\frac{\log p+10}{\pi\sqrt p},
\]
since $2+2\pi+\frac4\pi=9.55\ldots$ and $\frac1{3p^{3/2}}\le\frac{0.44}{\pi\sqrt p}$ for $p\ge3$.

\tools{Tome~4, Corollary~6.10 and Theorem~6.13 (Fejér, Weyl's criterion), Proposition~7.13 (Gauss sums).}
\leads Erdős and Turán (1948) proved the general form $D\le\frac C H+C\sum_{h\le H}\frac{|S(h)|}{hN}$ for every $H$; Vaaler (1985) found the optimal
smoothing functions, which give it with small explicit constants. For the squares modulo $p$, one finds numerically $D\sqrt p/\log p\approx0.14$
at $p=10007$, comfortably inside the bound of (c).
\end{solution}

\begin{solution}{pb:circlewalk}
\idea The Fourier coefficients of the law of $X_n$ are $n$-th powers $\cos^n(2\pi h\alpha)$. They tend to $0$ unless the cosine is $\pm1$, and
which frequencies survive is a question about $2h\alpha$ being an integer.

\writeup Since $e(hX_n)=e(h\alpha(\varepsilon_1+\cdots+\varepsilon_n))$ and the signs are independent,
\[
\mathbb E[e(hX_n)]=\prod_{k=1}^n\mathbb E[e(h\alpha\varepsilon_k)]=\cos^n(2\pi h\alpha),
\]
and $|\cos(2\pi h\alpha)|=1$ exactly when $2h\alpha\in\Z$.

(a) If $\alpha$ is irrational and $h\ne0$, then $2h\alpha\notin\Z$, so $\mathbb E[e(hX_n)]\to0=\int_0^1e(hx)\dd x$. By linearity $\mathbb E[P(X_n)]\to\int_0^1P$ for every
trigonometric polynomial $P$. Given a continuous $1$-periodic $g$ and $\varepsilon>0$, Fejér's theorem (Tome~4, Corollary~6.10(a)) gives a
trigonometric polynomial $P$ with $|g-P|\le\varepsilon$ everywhere, so $\limsup_n|\mathbb E[g(X_n)]-\int_0^1g|\le2\varepsilon$.

(b) Now $X_n$ lives in the finite set $\frac1q\Z/\Z$, and the law of a variable on this set is determined by the $q$ numbers $\mathbb E[e(hX)]$,
$0\le h<q$ (inversion of the DFT, Tome~4, Theorem~7.3); so it suffices to find the limits of these numbers. As $\gcd(a,q)=1$, $2h\alpha\in\Z$ iff
$q\mid2h$.

If $q$ is odd, this means $q\mid h$, and then $\cos(2\pi h\alpha)=1$; so $\mathbb E[e(hX_n)]\to1_{q\mid h}=\frac1q\sum_{j<q}e(\frac{hj}q)$, the coefficient of the uniform law on
$\frac1q\Z/\Z$.

If $q$ is even, $a$ is odd, and $q\mid2h$ means $h=l\frac q2$; then $\cos(2\pi h\alpha)=\cos(\pi la)=(-1)^l$. So $\mathbb E[e(hX_n)]$ tends to $0$ unless $h=l\frac q2$,
where it equals $(-1)^{ln}$. For $n$ even these are the coefficients of the uniform law on $\{0,\frac2q,\frac4q,\dots\}$, and for $n$ odd those of the
uniform law on $\{\frac1q,\frac3q,\dots\}$: indeed $\frac2q\sum_je(h(\frac1q+\frac{2j}q))=e(\frac hq)1_{\frac q2\mid h}=(-1)^l$ for $h=l\frac q2$. The two sets are disjoint, so
$X_n$ has no limit; this is the parity obstruction of Problem~\ref{pb:cyclewalk}, where $\alpha=\frac1k$.

\tools{Tome~4, Corollary~6.10 and Theorem~7.3; Tome~3, Chapter~5 (independence).}
\leads Compare with Weyl's theorem that $(\{n\alpha\})$ is equidistributed (Tome~4, Corollary~6.14): the deterministic walk and the random walk
are both governed by $\alpha$ being irrational. The \emph{speed} of convergence in (a), $\cos^n(2\pi h\alpha)\approx e^{-2\pi^2n\|2h\alpha\|^2}$ where $\|\cdot\|$ is the
distance to the nearest integer, depends on how well $2\alpha$ is approximated by rationals, that is, on its continued fraction (Tome~1,
Chapter~8).
\end{solution}

\begin{solution}{pb:kac}
\idea Moments of trigonometric sums count solutions of linear equations in the frequencies. With frequencies $2^k$, binary
expansions are unique and almost every solution cancels in pairs, exactly as for independent variables; that is why lacunary sums
behave like sums of independent terms. With frequencies $1,2,\dots,N$ there are far too many additive coincidences.

\writeup (a) Since $\cos A\cos B=\frac12(\cos(A-B)+\cos(A+B))$ and $\mathbb E[\cos mt]$ is $1$ for $m=0$ and $0$ for every other integer $m$, we get
$\mathbb E[\cos(at)\cos(bt)]=\frac12\cdot1_{a=b}$ for positive integers $a$, $b$. So both means vanish and both second moments are $\frac N2$.

(b) Expanding $\cos A_i=\frac12(e^{iA_i}+e^{-iA_i})$,
\[
\mathbb E\prod_{i=1}^4\cos(n_it)=\frac1{16}\#\Bigl\{\varepsilon\in\{\pm1\}^4:\sum_i\varepsilon_in_i=0\Bigr\},
\]
so $\mathbb E[S_N^4]=\frac M{16}$, where $M$ counts the $(k,\varepsilon)\in\{1,\dots,N\}^4\times\{\pm1\}^4$ with $\sum\varepsilon_i2^{k_i}=0$. In such a relation the terms with sign
$+$ and those with sign $-$ have equal sums; neither group is empty.

\emph{Two against two}, $2^a+2^b=2^c+2^d$. If $a\ne b$ the left side has exactly two binary digits $1$, at $a$ and $b$; the right side
then cannot be $2^{c+1}$, so $c\ne d$ and $\{c,d\}=\{a,b\}$. If $a=b$ the left side is $2^{a+1}$, which forces $c=d=a$. Either way the terms
cancel in pairs.

\emph{One against three}, $2^a=2^b+2^c+2^d$ with $b\le c\le d$. If $b,c,d$ are distinct the right side has three digits $1$; if they are
equal it is $3\cdot2^b$; if $b<c=d$ it is $2^b+2^{c+1}$ with two digits. So $b=c<d$, and $2^{b+1}+2^d$ is a power of $2$ only if $d=b+1$: the
relation is $2^{m+2}=2^{m+1}+2^m+2^m$, with $1\le m\le N-2$.

Counting. Pairings: choose one of the $3$ ways of splitting the four positions into two pairs, an exponent for each pair ($N^2$
ways) and the sign of the first element of each pair ($4$ ways): $12N^2$. A tuple is counted more than once only if it fits two
splittings, which happens exactly when the four exponents are equal (two signs $+$, two signs $-$: $6N$ tuples, each fitting $2$
splittings); so there are $12N^2-6N$ pairing tuples. The other relations: $N-2$ values of $m$, $\frac{4!}{2!}=12$ arrangements of the exponents
$m+2,m+1,m,m$, and $2$ sign choices: $24(N-2)$ tuples, none of which cancels in pairs. So
\[
M=12N^2-6N+24(N-2)=12N^2+18N-48,\qquad\mathbb E[S_N^4]=\frac34N^2+\frac98N-3,
\]
and $\mathbb E[S_N^4]/\mathbb E[S_N^2]^2=3+\frac9{2N}-\frac{12}{N^2}\to3$.

(c) Now count the solutions of $\sum\varepsilon_ik_i=0$ with $1\le k_i\le N$. For each of the $6$ sign patterns with two $+$ and two $-$, the number of
solutions of $k_1+k_2=k_3+k_4$ is $\sum_{s=2}^{2N}r(s)^2$ with $r(s)=\min(s-1,2N+1-s)$, that is $2\sum_{j=1}^{N-1}j^2+N^2=\frac{2N^3+N}3$. For each of the $8$
patterns with three signs alike, the number of solutions of $k_1+k_2+k_3=k_4$ is $\sum_{s\le N}\binom{s-1}2=\binom N3$. So
\[
\mathbb E[T_N^4]=\frac1{16}\Bigl(2(2N^3+N)+8\binom N3\Bigr)\sim\frac1{16}\Bigl(4+\frac43\Bigr)N^3=\frac{N^3}3,
\]
and $\mathbb E[T_N^4]/\mathbb E[T_N^2]^2\sim\frac{4N}3\to\infty$.

\tools{Tome~0, Chapter~2 (Euler's formula); Tome~1, Chapter~1 (binary expansions); Tome~3, Chapter~4 (moments).}
\leads The same counting shows that every moment of $S_N/\sqrt{N/2}$ tends to the corresponding Gaussian moment, and the method of
moments then gives Kac's theorem (1946): $S_N/\sqrt{N/2}$ converges in law to $N(0,1)$, although the terms $\cos(2^kt)$ are not
independent. Salem and Zygmund (1947) extended it to $\sum\cos(n_kt)$ whenever $n_{k+1}/n_k\ge q>1$. Lacunarity alone is not enough for
other functions: Erdős and Fortet observed that for $f(t)=\cos t+\cos2t$ and $n_k=2^k-1$, the limit law of $\sum_kf(n_kt)/\sqrt N$ is a mixture of
Gaussians, because of the relations $2n_k+1=n_{k+1}$.
\end{solution}

\begin{solution}{pb:cantor}
\idea A sum of independent digits has a product characteristic function. Multiplying $t$ by $3$ shifts the product by one factor,
so along $t=2\pi3^m$ the function repeats itself and cannot decay; a density would force decay.

\writeup (a) Let $X_K=\sum_{k\le K}2B_k3^{-k}$. By independence,
\[
\mathbb E[e^{itX_K}]=\prod_{k\le K}\frac{1+e^{2it/3^k}}2=\prod_{k\le K}e^{it/3^k}\cos\frac t{3^k}=e^{it(1-3^{-K})/2}\prod_{k\le K}\cos\frac t{3^k}.
\]
As $|X-X_K|\le3^{-K}$, $|e^{itX}-e^{itX_K}|\le|t|3^{-K}$, so the left side tends to $\varphi(t)$; hence the product converges and the formula holds.

(b) For $t=2\pi3^m$, the factors with $k\le m$ are $\cos(2\pi3^{m-k})=1$, and the others are $\cos(2\pi3^{-j})$ with $j=k-m\ge1$. So
$|\varphi(2\pi3^m)|=\bigl|\prod_{j\ge1}\cos\frac{2\pi}{3^j}\bigr|=|\varphi(2\pi)|$. The first factor is $-\frac12$. For $j\ge2$, $x_j=\frac{2\pi}{3^j}\le\frac{2\pi}9<\frac\pi4$, and for
$|x|\le\frac\pi4$, $\cos x\ge1-\frac{x^2}2\ge e^{-x^2}$ (as $1-u\ge e^{-2u}$ for $0\le u\le\frac12$). So $\prod_{j\ge2}\cos x_j\ge\exp(-\sum_{j\ge2}x_j^2)=e^{-\pi^2/18}>0$,
and $|\varphi(2\pi)|>0$; numerically it is $0.3714\ldots$

If $X$ had a density $f$, then $\varphi(t)=\int f(x)e^{itx}\dd x\to0$ as $t\to\infty$ (Riemann--Lebesgue): for $f=1_{[a,b]}$, $|\int_a^be^{itx}\dd x|\le\frac2{|t|}$; this
extends to step functions by linearity, and to every integrable $f$ because step functions approximate $f$ in mean,
$\int|f-s|\le\varepsilon$, and $|\int(f-s)e^{itx}|\le\varepsilon$. But $\varphi(2\pi3^m)$ does not tend to $0$. So $X$ has no density.

(c) From $\sin2u=2\sin u\cos u$, by induction
\[
\sin t=2^m\sin\frac t{2^m}\prod_{k=1}^m\cos\frac t{2^k},
\]
and $2^m\sin\frac t{2^m}\to t$. So $\prod_{k\ge1}\cos\frac t{2^k}=\frac{\sin t}t$ for $t\ne0$. As in (a), $Y=\sum_k\varepsilon_k2^{-k}$ has characteristic function
\[
\prod_k\cos\frac t{2^k}=\frac{\sin t}t=\frac12\int_{-1}^1e^{itu}\dd u,
\]
that of the
uniform law on $[-1,1]$; by uniqueness (Tome~3, Theorem~6.8), $Y$ is uniform.

\tools{Tome~3, Exercise~5.7 and Theorem~6.8 (the Cantor law, characteristic functions); Tome~0, Chapter~2 (trigonometry).}
\leads The same question for $Y_\lambda=\sum_k\varepsilon_k\lambda^k$, the \emph{Bernoulli convolutions}, is a famous problem. For $\lambda<\frac12$ the law lives on a
Cantor set; for $\lambda=\frac12$ it is uniform. For $\frac12<\lambda<1$, Erdős (1939) used exactly the argument of (b) to prove that the law is
singular when $\frac1\lambda$ is a Pisot number, such as the golden ratio: then $\theta^m$ is very close to an integer, $\theta=\frac1\lambda$, so the
characteristic function does not decay along $t=2\pi\theta^m$. Solomyak (1995) proved that the law has a density for almost every
$\lambda\in\left]\frac12,1\right[$; which $\lambda$ exactly is still open.
\end{solution}

\begin{solution}{pb:sincsum}
\idea Poisson summation turns $\sum_nf(n)$ into $\sum_m\hat f(m)$, and $\hat f(0)$ is the integral. For $\operatorname{sinc}^k$ the transform is a probability
density supported on $[-\frac k{2\pi},\frac k{2\pi}]$, so all the terms $m\ne0$ vanish as long as $k<2\pi$. The first integer past $2\pi$ is $7$.

\writeup (a) The $2\pi$-periodic function equal to $\frac{\pi-x}2$ on $]0,2\pi[$ is piecewise $C^1$, with $c_0=0$ and, integrating by parts,
$c_n=\frac1{2\pi}\int_0^{2\pi}\frac{\pi-x}2e^{-inx}\dd x=\frac1{2in}$ for $n\ne0$. Its Fourier partial sums are $\sum_{n=1}^N\frac{\sin nx}n$, and Dirichlet's theorem
(Tome~4, Theorem~6.5) at the point of continuity $x=1$ gives $\sum_{n\ge1}\frac{\sin n}n=\frac{\pi-1}2$. Hence
$\sum_{n\in\Z}\operatorname{sinc}n=1+2\cdot\frac{\pi-1}2=\pi$, the value of Dirichlet's integral $\int_{\R}\operatorname{sinc}$ (Tome~4, Chapter~4).

(b) $\operatorname{sinc}x=\frac12\int_{-1}^1e^{ixu}\dd u=\mathbb E[e^{ixU}]$, so by independence $f(x)=\mathbb E[e^{ixS}]=\int g_k(s)e^{ixs}\dd s$, where $S=U_1+\cdots+U_k$ takes values in
$[-k,k]$; for $k\ge2$, $g_k$ is continuous. Put $G(\eta)=2\pi g_k(2\pi\eta)$, continuous, integrable and supported on $[-\frac k{2\pi},\frac k{2\pi}]$. With $s=2\pi\eta$,
$\hat G(x)=\int g_k(s)e^{-ixs}\dd s=f(-x)=f(x)$, which is integrable for $k\ge2$. By Fourier inversion (Tome~4, Theorem~8.6),
$G(\eta)=\int f(x)e^{2\pi ix\eta}\dd x=\hat f(-\eta)=\hat f(\eta)$, as $f$ is even.

(c) For $k\ge2$, $f$ is continuous with $|f(x)|\le\min(1,x^{-2})\le\frac4{(1+|x|)^2}$, and $\hat f$ is bounded with compact support, so Poisson summation
(Tome~4, Theorem~8.9) applies:
\[
\sum_{n\in\Z}\operatorname{sinc}^kn=\sum_{m\in\Z}\hat f(m)=\int_{\R}\operatorname{sinc}^k+2\pi\sum_{m\ne0}g_k(2\pi m).
\]
For $k\le6$, $g_k$ vanishes outside $[-6,6]$ and $2\pi>6$, so the extra terms vanish. For $k=7$, only $m=\pm1$ survive ($4\pi>7$), and $2\pi\in[5,7]$.
Writing $U_i=1-2V_i$ with $V_i$ uniform on $[0,1]$, $S=7-2(V_1+\cdots+V_7)$, and for $0\le y\le1$, $P(V_1+\cdots+V_7\le y)=\frac{y^7}{7!}$ (the volume of a
simplex), so $V_1+\cdots+V_7$ has density $\frac{y^6}{6!}$ there. Hence $g_7(s)=\frac12\cdot\frac{((7-s)/2)^6}{6!}=\frac{(7-s)^6}{2^7\cdot6!}$ for $5\le s\le7$, and the excess is
$2\cdot2\pi g_7(2\pi)=\frac{4\pi(7-2\pi)^6}{2^7\cdot6!}=1.849\ldots\cdot10^{-5}$ (Figure~\ref{fig:sinc}).

\tools{Tome~4, Theorems~6.5, 8.6 and 8.9 (Dirichlet, Fourier inversion, Poisson summation); Tome~3, Chapter~6 (characteristic functions).}
\leads Baillie, Borwein and Borwein (2008) found the general rule:
\[
\sum_{n\in\Z}\prod_j\operatorname{sinc}(a_jn)=\int_{\R}\prod_j\operatorname{sinc}(a_jx)\dd x\qquad\text{whenever }\sum_ja_j<2\pi,
\]
by the same argument. The mechanism is the one behind the Borwein integrals of Problem~\ref{pb:borwein}: a density whose support grows with
the number of factors, and an identity that holds until the support crosses a threshold ($2a_0$ there, $2\pi$ here). In signal processing this is aliasing: sampling at the integers is faithful to frequencies below the
Nyquist limit (Tome~4, Chapter~8).
\end{solution}

\begin{solution}{pb:influence}
\idea Flipping voter $i$ matters exactly where the discrete derivative in the $i$-th direction is non-zero, and the derivative of a
multilinear polynomial is read off its coefficients. Parseval then turns probabilities into sums of squares.

\writeup (a) Let $D_if(x)=\frac12\bigl(f(x^{i\to1})-f(x^{i\to-1})\bigr)$, where $x^{i\to\pm1}$ is $x$ with $x_i$ set to $\pm1$. On the expansion, $D_ix^S=x^{S\setminus\{i\}}$ if
$i\in S$ and $0$ otherwise, so $D_if=\sum_{S\ni i}\hat f(S)x^{S\setminus\{i\}}$. For Boolean $f$, $D_if$ takes values in $\{-1,0,1\}$, and $|D_if(x)|=1$ exactly when
$f(x)\ne f(x^{\oplus i})$. Hence, by Parseval (the monomials $x^{S\setminus\{i\}}$, $S\ni i$, are distinct),
\[
\operatorname{Inf}_i(f)=\mathbb E[(D_if)^2]=\sum_{S\ni i}\hat f(S)^2,\qquad I(f)=\sum_i\sum_{S\ni i}\hat f(S)^2=\sum_S|S|\hat f(S)^2.
\]
And $\operatorname{Var}f=\sum_{S\ne\emptyset}\hat f(S)^2\le\sum_S|S|\hat f(S)^2$, with equality iff $\hat f(S)=0$ whenever $|S|\ge2$.

(b) If $f$ is monotone, $D_if\in\{0,1\}$, so $(D_if)^2=D_if$ and $\operatorname{Inf}_i(f)=\mathbb E[D_if]=\hat f(\{i\})$, the constant term of $D_if$. By Cauchy--Schwarz and
Parseval, $I(f)=\sum_i\hat f(\{i\})\le\sqrt n\bigl(\sum_i\hat f(\{i\})^2\bigr)^{1/2}\le\sqrt n$.

(c) Voter $i$ is pivotal for majority iff the other $n-1$ votes split evenly, which has probability $\binom{n-1}{(n-1)/2}2^{1-n}$. With $n-1=2m$ and
$\binom{2m}m\sim\frac{4^m}{\sqrt{\pi m}}$ (Stirling), $\operatorname{Inf}_i\sim\frac1{\sqrt{\pi m}}\sim\sqrt{\frac2{\pi n}}$, so $I(\operatorname{Maj}_n)=n\operatorname{Inf}_1\sim\sqrt{\frac{2n}\pi}$. Majority is monotone, so
$\hat f(\{i\})=\operatorname{Inf}_i$ by (b), and $W^1(\operatorname{Maj}_n)=n\operatorname{Inf}_1^2\to\frac2\pi$.

\tools{Tome~3, Chapter~2 (binomial law) and Tome~0, Chapter~5 (Stirling's formula); Tome~0, Chapter~1 (Cauchy--Schwarz).}
\leads So among monotone rules, majority gives the voters almost the most power possible: $\sqrt{2/\pi}=0.80$ of the bound $\sqrt n$. At the
other extreme, Kahn, Kalai and Linial (1988) proved that every Boolean $f$ with $\operatorname{Var}f=1$ has a voter with influence at least $c\frac{\log n}n$,
and the ``tribes'' function of Ben-Or and Linial shows that this is sharp; the proof uses the Bonami lemma of Problem~\ref{pb:bonami}.
For $f=1-2\cdot1_A$, $I(f)$ is the number of edges of the cube leaving $A$ divided by $2^{n-1}$, which links (a) with the edge-isoperimetric
inequality of Problem~\ref{pb:cubeiso}.
\end{solution}

\begin{solution}{pb:granularity}
\idea Boolean functions are integer-valued, so in $\{0,1\}$ coordinates their polynomial has integer coefficients; changing back to
$\pm1$ coordinates costs a factor $2$ per variable of a monomial, at most $2^d$. A coarse spectrum cannot be spread thin, and the
same integrality shows that every relevant variable carries a definite amount of influence.

\writeup (a) Let $g=\frac{1-f}2\in\{0,1\}$ and $y_i=\frac{1-x_i}2\in\{0,1\}$. As a function of $y\in\{0,1\}^n$, $g$ has a unique multilinear expansion
$g=\sum_Ta_Ty^T$, and evaluating at $y=1_V$ gives $g(1_V)=\sum_{T\subseteq V}a_T$; by Möbius inversion on subsets, $a_T=\sum_{U\subseteq T}(-1)^{|T\setminus U|}g(1_U)$,
an integer. The substitutions $x_i=1-2y_i$ and $y_i=\frac{1-x_i}2$ are affine in each variable and preserve multilinearity, so they do not
raise the degree in either direction; as expansions are unique, $\deg_yg=\deg_xg=d$, and $a_T=0$ for $|T|>d$. Now
$y^T=2^{-|T|}\prod_{i\in T}(1-x_i)$ has $x$-coefficients in $2^{-|T|}\Z\subseteq2^{-d}\Z$, so every $\hat g(S)$ is in $2^{-d}\Z$, and $\hat f(S)=1_{S=\emptyset}-2\hat g(S)\in2^{1-d}\Z$.

(b) Each non-zero $\hat f(S)^2$ is at least $4^{1-d}$, and they add up to $1$.

(c) Induction on $n$. For $n=0$, $p$ is a non-zero constant. Write $p=x_nq+r$, with $q$, $r$ multilinear in $x'=(x_1,\dots,x_{n-1})$ and
$\deg q\le k-1$. If $q=0$, then $p=r\ne0$ and the induction hypothesis applies. Otherwise, for each $x'$ with $q(x')\ne0$, at least one
of $p(x',1)=r+q$ and $p(x',-1)=r-q$ is non-zero; so $P(p\ne0)\ge\frac12P(q\ne0)\ge\frac12\cdot2^{-(k-1)}=2^{-k}$.

(d) $D_if=\sum_{S\ni i}\hat f(S)x^{S\setminus\{i\}}$ (Problem~\ref{pb:influence}) has degree at most $d-1$, is non-zero iff $i$ is relevant, and
$\operatorname{Inf}_i(f)=P(D_if\ne0)$. By (c), $\operatorname{Inf}_i(f)\ge2^{1-d}$ for relevant $i$. Since $I(f)=\sum_S|S|\hat f(S)^2\le d$, the number $R$ of relevant variables
satisfies $R\,2^{1-d}\le d$, that is $R\le d\,2^{d-1}$.

\tools{Tome~3, Chapter~1 (inclusion--exclusion, Möbius inversion on subsets); Tome~0, Chapter~3 (polynomials).}
\leads The bound $d\,2^{d-1}$ is due to Nisan and Szegedy (1994), in their study of the degree of Boolean functions as real
polynomials. Chiarelli, Hatami and Saks (2020) improved it to $O(2^d)$, which is the right order: the ``address'' function, where $d-1$
bits point to one of $2^{d-1}$ other bits and the output is the bit pointed to, has degree $d$ and $d-1+2^{d-1}$ relevant variables.
\end{solution}

\begin{solution}{pb:bent}
\idea Parseval forces the $2^n$ squared coefficients to average $2^{-n}$, so the largest is at least that; and $2^n\hat f(a)$ is an integer,
which decides when they can all be equal. Distance to linear functions is correlation in disguise.

\writeup (a) $\sum_a\hat f(a)^2=\mathbb E[f^2]=1$ over $2^n$ values of $a$, so some $\hat f(a)^2\ge2^{-n}$. If $f$ is bent, $2^n\hat f(a)$, a sum of $2^n$ signs, is an
integer equal to $\pm2^{n/2}$, and $2^{n/2}$ is an integer only for $n$ even.

(b) $\hat f(a,b)=2^{-2m}\sum_x(-1)^{a\cdot x}\sum_y(-1)^{(x+b)\cdot y}$. The inner sum is $2^m$ if $x=b$ and $0$ otherwise, so $\hat f(a,b)=2^{-m}(-1)^{a\cdot b}$, of modulus
$2^{-m}=2^{-n/2}$ with $n=2m$.

(c) For $g=\pm(-1)^{a\cdot x}$, the number of disagreements is $\frac12\sum_x(1-f(x)g(x))=2^{n-1}(1\mp\hat f(a))$. Minimising over $a$ and the sign gives
$2^{n-1}(1-\max_a|\hat f(a)|)$, which is at most $2^{n-1}(1-2^{-n/2})$ by (a). Equality means $\max_a|\hat f(a)|=2^{-n/2}$; as the squares average $2^{-n}$,
they are then all equal to $2^{-n}$, and $f$ is bent.

(d) $\hat f(0)=2^{-n}\sum_xf(x)=\pm2^{-n/2}$, so $\sum_xf(x)=\pm2^{n/2}$, and the number of $x$ with $f(x)=-1$ is $\frac12(2^n\mp2^{n/2})=2^{n-1}\mp2^{n/2-1}$.

\tools{Tome~2, Chapter~11 (characters of finite abelian groups); Tome~4, Chapter~7 (Parseval).}
\leads Bent functions were defined and named by Oscar Rothaus in the 1960s, in research that was published only in 1976. A cipher whose
components are close to linear functions is vulnerable to linear cryptanalysis (Matsui, 1993), so high nonlinearity is a design
goal; but by (d) bent functions are unbalanced, and cryptographers trade a little nonlinearity for balance. The set
$\{x:f(x)=-1\}$ of a bent function is a \emph{Hadamard difference set}, and its $\pm1$ matrix $(f(x+y))_{x,y}$ is a Hadamard matrix
(Problem~\ref{pb:hadamard}).
\end{solution}

\begin{solution}{pb:majoritystable}
\idea Noise stability weights the spectrum by $\rho^{|S|}$, so it is a generating function of the spectral distribution. For majority,
the central limit theorem turns the two vote counts into a Gaussian pair, and the probability that two correlated Gaussians have
the same sign is an angle.

\writeup (a) $\mathbb E[x^Sy^T]=\prod_i\mathbb E[x_i^{1_S(i)}y_i^{1_T(i)}]$ by independence of the pairs; a factor with $i$ in exactly one of $S$, $T$ is $0$, and a
factor with $i\in S\cap T$ is $\rho$. So $\mathbb E[x^Sy^T]=\rho^{|S|}1_{S=T}$, and expanding $f(x)f(y)$ gives $\operatorname{Stab}_\rho(f)=\sum_S\rho^{|S|}\hat f(S)^2$.

(b) Let $X$, $Z$ be independent $N(0,1)$ and $Y=\rho X+\sqrt{1-\rho^2}Z$; then $(X,Y)$ is a Gaussian pair with the right covariance, so it has the right
law (Tome~3, Theorem~6.15). For a rotation $A$ of the plane, $A(X,Z)$ is Gaussian with covariance $AA^{\mathsf T}=I$, so it has the same law as
$(X,Z)$. Rotations permute the open sectors with vertex $0$ and a given angle, so the probability of such a sector depends only on its
angle $\theta$; it is additive in $\theta$, and the full plane has probability $1$, so it equals $\frac\theta{2\pi}$ for rational multiples of $2\pi$, and by
monotonicity for all $\theta$. Write $\rho=\cos\beta$ with $0<\beta<\pi$ and $(X,Z)=R(\cos\Theta,\sin\Theta)$; then $Y=R\cos(\Theta-\beta)$, and $\{X>0,Y>0\}$ is the sector
$\{|\Theta|<\frac\pi2,\ |\Theta-\beta|<\frac\pi2\}$ of angle $\pi-\beta$. So
\[
P(X>0,Y>0)=\frac{\pi-\arccos\rho}{2\pi}=\frac14+\frac{\arcsin\rho}{2\pi},
\]
and $\mathbb E[\operatorname{sgn}X\operatorname{sgn}Y]=2P(\operatorname{sgn}X=\operatorname{sgn}Y)-1=4P(X>0,Y>0)-1=\frac2\pi\arcsin\rho$.

(c) The vectors $(x_i,y_i)$ are independent and identically distributed, with mean $0$ and covariance matrix $\begin{pmatrix}1&\rho\\\rho&1\end{pmatrix}$, so
$\frac1{\sqrt n}\sum_i(x_i,y_i)$ converges in law to $(X,Y)$. For $n$ odd the sums are never $0$, and $\operatorname{Stab}_\rho(\operatorname{Maj}_n)=4P(\sum x_i>0,\sum y_i>0)-1$ by the
symmetry $(x,y)\mapsto(-x,-y)$. The quadrant $\{u>0,v>0\}$ has a boundary of probability $0$ for $(X,Y)$, so its probability converges, and
$\operatorname{Stab}_\rho(\operatorname{Maj}_n)\to\frac2\pi\arcsin\rho$.

(d) Let $G_n(\rho)=\operatorname{Stab}_\rho(\operatorname{Maj}_n)=\sum_kW^k_n\rho^k$, where $W^k_n=W^k(\operatorname{Maj}_n)\ge0$ and $\sum_kW^k_n=1$, and $G(\rho)=\frac2\pi\arcsin\rho=\sum_kb_k\rho^k$ with
$b_{2j+1}=\frac2\pi\binom{2j}j\frac1{4^j(2j+1)}$ and $b_{2j}=0$ (integrate $(1-u^2)^{-1/2}=\sum_j\binom{2j}j\frac{u^{2j}}{4^j}$). Take any subsequence; by a diagonal extraction, a
further subsequence has $W^k_n\to w_k$ for every $k$. For $0\le\rho<1$, the tail $\sum_{k\ge K}W^k_n\rho^k$ is at most $\frac{\rho^K}{1-\rho}$, uniformly in $n$, so
$\sum_kw_k\rho^k=G(\rho)$ on $[0,1[$. Two power series that agree on $[0,1[$ have the same coefficients (let $\rho\to0$, divide by $\rho$, repeat), so
$w_k=b_k$. Every subsequence has a further subsequence with these limits, so $W^k_n\to b_k$. In particular $W^1\to b_1=\frac2\pi$, as in
Problem~\ref{pb:influence}, and $W^3\to\frac1{3\pi}$.

\tools{Tome~3, Theorem~6.15 and the remark after it (Gaussian vectors, the CLT in $\R^2$); Tome~4, Chapter~3 (power series).}
\leads Sheppard found his formula in 1899, studying correlations in biometric data. The limit $\frac2\pi\arcsin\rho$ is the starting point of
the \emph{Majority is Stablest} theorem of Mossel, O'Donnell and Oleszkiewicz (2010): among Boolean functions with mean $0$ and all
influences small, none is more noise-stable than majority. With the unique games conjecture, it implies that the Goemans--Williamson
algorithm for maximum cut, whose approximation ratio $0.878\ldots$ comes from the same arcsine, is optimal.
\end{solution}

\begin{solution}{pb:arrow}
\idea One voter's three pairwise preferences are three correlated signs, and ``not all equal'' is a quadratic polynomial in three
signs. Taking expectations gives noise stability at the correlation $-\frac13$, and the sign of $(-\frac13)^k$ separates level $1$ from all others.

\writeup (a) The six rankings $a{>}b{>}c$, $a{>}c{>}b$, $b{>}a{>}c$, $b{>}c{>}a$, $c{>}a{>}b$, $c{>}b{>}a$ give $(x_i,y_i,z_i)=(1,1,-1)$, $(1,-1,-1)$, $(-1,1,-1)$, $(-1,1,1)$,
$(1,-1,1)$, $(-1,-1,1)$: the six vectors other than $\pm(1,1,1)$, each once (a cycle would need a cyclic individual ranking). Each coordinate is
uniform, and $x_iy_i=1$ for exactly two of the six, so $\mathbb E[x_iy_i]=\frac{2-4}6=-\frac13$; the same holds for $(y_i,z_i)$ and $(z_i,x_i)$ by the cyclic
symmetry $a\to b\to c\to a$. Voters are independent, so the three pairs are $(-\frac13)$-correlated.

(b) For $u,v,w\in\{-1,1\}$: if $u=v=w$, $uv+vw+wu=3$; otherwise exactly one product is $1$ and $uv+vw+wu=-1$. So
$1-1_{u=v=w}=\frac34-\frac14(uv+vw+wu)$. With $u=f(x)$, $v=f(y)$, $w=f(z)$ and (a),
\[
P(\text{rational})=\frac34-\frac14\bigl(\mathbb E[f(x)f(y)]+\mathbb E[f(y)f(z)]+\mathbb E[f(z)f(x)]\bigr)=\frac34-\frac34\operatorname{Stab}_{-1/3}(f).
\]

(c) $\operatorname{Maj}_3=\frac12(x_1+x_2+x_3)-\frac12x_1x_2x_3$ (check the eight points), so $W^1=\frac34$, $W^3=\frac14$, $\operatorname{Stab}_{-1/3}=-\frac14-\frac1{108}=-\frac7{27}$, and
$P=\frac34+\frac7{36}=\frac{17}{18}$: a cycle has probability $\frac1{18}$, that is $12$ of the $216$ profiles, the answer to the think box. Majority is odd,
$f(-x)=-f(x)$, so only odd $|S|$ occur and $\operatorname{Stab}_{-\rho}=-\operatorname{Stab}_\rho$. By Problem~\ref{pb:majoritystable}, $\operatorname{Stab}_{-1/3}(\operatorname{Maj}_n)\to-\frac2\pi\arcsin\frac13$, and
$P\to\frac34+\frac3{2\pi}\arcsin\frac13=0.91226\ldots$; a cycle has limiting probability $0.0877\ldots$ (Figure~\ref{fig:condorcet}).

(d) If every profile is rational, $P=1$ and $\operatorname{Stab}_{-1/3}(f)=-\frac13$, that is $\sum_k\bigl((-\frac13)^k+\frac13\bigr)W^k(f)=0$ since $\sum_kW^k=1$. Each bracket is
$\ge0$, and is $0$ only for $k=1$; so all the spectral weight is on level $1$, and $f=\sum_ia_ix_i$. Flipping $x_i$ changes $f$ by $2a_ix_i\in\{0,\pm2\}$, so
$a_i\in\{0,\pm1\}$, and $\sum a_i^2=1$ leaves one non-zero $a_i$: $f=\pm x_i$. Unanimity excludes $-x_i$.

\tools{Problems~\ref{pb:majoritystable} and~\ref{pb:influence}; Tome~3, Chapter~2 (expectation).}
\leads Arrow's theorem in general (any number of candidates and any rules for the different pairs, under the conditions of
unanimity and independence of irrelevant alternatives) needs a combinatorial proof; Kalai's formula gives the case above and
more importantly, a robust version: if $P(\text{rational})\ge1-\varepsilon$, then by the Friedgut--Kalai--Naor theorem $f$ is $O(\varepsilon)$-close to a
dictatorship or an anti-dictatorship. Mossel (2012) extended this quantitative Arrow theorem to any number of candidates and to different rules for different
pairs.
\end{solution}

\begin{solution}{pb:bonami}
\idea Peel off one variable: $f=x_ng+h$. The fourth moment splits into three terms, and the cross term is controlled by Cauchy--Schwarz
and induction; the constants are exactly what makes $9^d$ reproduce itself.

\writeup (a) Induction on $n$; for $n=0$, $f$ is a constant $c$ and $\mathbb E[f^4]=c^4=\mathbb E[f^2]^2$. Write $f=x_ng+h$ with $g$, $h$ multilinear in $x_1,\dots,x_{n-1}$,
$\deg g\le d-1$ and $\deg h\le d$ (if $d=0$, $g=0$). As $x_n$ is independent of $g$, $h$ with $\mathbb E[x_n]=\mathbb E[x_n^3]=0$ and $x_n^2=1$,
\[
\mathbb E[f^2]=\mathbb E[g^2]+\mathbb E[h^2],\qquad\mathbb E[f^4]=\mathbb E[h^4]+6\,\mathbb E[g^2h^2]+\mathbb E[g^4].
\]
By induction $\mathbb E[h^4]\le9^d\mathbb E[h^2]^2$ and $\mathbb E[g^4]\le9^{d-1}\mathbb E[g^2]^2$, and by Cauchy--Schwarz $\mathbb E[g^2h^2]\le\sqrt{\mathbb E[g^4]\mathbb E[h^4]}\le3^{2d-1}\mathbb E[g^2]\mathbb E[h^2]$.
Since $6\cdot3^{2d-1}=2\cdot9^d$,
\[
\mathbb E[f^4]\le9^d\bigl(\mathbb E[h^2]^2+2\mathbb E[g^2]\mathbb E[h^2]+\mathbb E[g^2]^2\bigr)=9^d\,\mathbb E[f^2]^2.
\]

(b) Let $s=\mathbb E[f^2]>0$ and $q=P(f^2\ge\frac s4)$. Then $s\le\frac s4+\mathbb E[f^21_{f^2\ge s/4}]\le\frac s4+\sqrt{\mathbb E[f^4]}\sqrt q$, so
$q\ge\frac{9s^2}{16\,\mathbb E[f^4]}\ge\frac9{16}9^{-d}$ by (a).

(c) Split $md$ variables into $d$ blocks $B_1,\dots,B_d$ of size $m$, let $L_j=m^{-1/2}\sum_{i\in B_j}x_i$, and $f=L_1\cdots L_d$, multilinear of degree $d$. The
$L_j$ are independent, $\mathbb E[L_j^2]=1$ and $\mathbb E[L_j^4]=\frac1{m^2}\bigl(m+3m(m-1)\bigr)=3-\frac2m$. So $\mathbb E[f^2]=1$ and $\mathbb E[f^4]=(3-\frac2m)^d>C^d$ once $\frac2m<3-C$.

\tools{Tome~3, Chapter~4 (Cauchy--Schwarz for expectations, moments of sums of signs).}
\leads Bonami proved this in 1970, in a study of Fourier coefficients on groups; Beckner and Gross (1975) found the sharp form,
\emph{hypercontractivity}: $\|T_\rho f\|_4\le\|f\|_2$ for $\rho\le\frac1{\sqrt3}$, where $T_\rho$ multiplies $\hat f(S)$ by $\rho^{|S|}$. Applied to the derivatives
$D_if$, it is the key to the theorem of Kahn, Kalai and Linial quoted in Problem~\ref{pb:influence}, to Friedgut's theorem that functions
of small total influence depend essentially on few variables, and to much of the analysis of Boolean functions in complexity
theory. Part~(c) only shows that the base is at least $3$, but $9$ is in fact the right base: by the central limit theorem, polynomials on
the cube approximate polynomials in Gaussian variables, and for the Hermite polynomials $H_d$ the $d$-th root of
$\mathbb E[H_d^4]/\mathbb E[H_d^2]^2$ tends to $9$ (slowly: one finds numerically $7.8$ at $d=30$).
\end{solution}

\begin{solution}{pb:hypercubemix}
\idea The characters of the cube diagonalise the walk: $x^S$ decays at rate $1-\frac{|S|}n$. Cauchy--Schwarz turns the total variation distance
into a sum over all characters, dominated by the $n$ slowest ones; one of them, $\sum x_i$, also gives the lower bound.

\writeup (a) At step $t+1$, a coordinate $I$ is chosen and given a fresh sign $\xi$. If $I\notin S$ (probability $1-\frac{|S|}n$), $X^S$ is unchanged; if
$I\in S$, $X_{t+1}^S=\xi X_t^{S\setminus\{I\}}$ has conditional mean $0$. So $\mathbb E[X_{t+1}^S]=(1-\frac{|S|}n)\mathbb E[X_t^S]$, and $X_0^S=1$.

(b) Let $\pi_t(x)=P(X_t=x)$ and $h_t=2^n\pi_t$, so that $\mathbb E_U[h_t]=1$ under the uniform law $U$, and $\hat h_t(S)=\mathbb E_U[h_tx^S]=\sum_x\pi_t(x)x^S=\mathbb E[X_t^S]$.
Then
\[
d_n(t)=\frac12\sum_x\bigl|\pi_t(x)-2^{-n}\bigr|=\frac12\mathbb E_U|h_t-1|\le\frac12\mathbb E_U\bigl[(h_t-1)^2\bigr]^{1/2}=\frac12\Bigl(\sum_{S\ne\emptyset}\hat h_t(S)^2\Bigr)^{1/2},
\]
which is the first bound by (a). With $\binom nk\le\frac{n^k}{k!}$ and $(1-\frac kn)^{2t}\le e^{-2tk/n}$, the sum is at most $\sum_{k\ge1}\frac{(ne^{-2t/n})^k}{k!}=\exp(ne^{-2t/n})-1$.

(c) By (a), $\mu_t=\mathbb E[W_t]=n(1-\frac1n)^t$ and $\mathbb E[W_t^2]=n+n(n-1)(1-\frac2n)^t$, so
\[
\operatorname{Var}W_t=n+n(n-1)\Bigl(1-\frac2n\Bigr)^t-n^2\Bigl(1-\frac1n\Bigr)^{2t}\le n,
\]
because $1-\frac2n\le(1-\frac1n)^2$. Under $U$, $W$ has mean $0$ and variance $n$. Let $A=\{W>\frac{\mu_t}2\}$. By Chebyshev, $U(A)\le\frac{4n}{\mu_t^2}$ and
$P(X_t\notin A)\le P(|W_t-\mu_t|\ge\frac{\mu_t}2)\le\frac{4n}{\mu_t^2}$, so $d_n(t)\ge P(X_t\in A)-U(A)\ge1-\frac{8n}{\mu_t^2}$.

(d) For $t\ge\frac12n\log n+cn$, $ne^{-2t/n}\le e^{-2c}$, and (b) gives the upper bound. For $t=\lfloor\frac12n\log n-cn\rfloor$, use $\log(1-u)\ge-u-u^2$ for
$0\le u\le\frac12$: $\mu_t\ge ne^{-t/n-t/n^2}\ge\sqrt n\,e^ce^{-\log n/(2n)}$, so $\frac{8n}{\mu_t^2}\le8e^{-2c}e^{\log n/n}\to8e^{-2c}$, and (c) gives the lower bound.

\tools{Tome~3, Chapter~8 (total variation distance, mixing) and Chapter~4 (Chebyshev's inequality).}
\leads The coupon collector needs about $n\log n$ steps to refresh every coordinate, which gives a cheap upper bound twice too large;
the Fourier analysis sees that the walk is mixed after half of that, when about $\sqrt n$ coordinates are still untouched, too few to be
detected. Aldous and Diaconis (1986) named the cutoff phenomenon, after Diaconis and Shahshahani had found it for random
transpositions; on the cycle of Tome~3, Theorem~8.17, there is no cutoff, the distance decaying gradually on the scale $n^2$.
\end{solution}

\begin{solution}{pb:fejer}
\idea The polynomial $z^nT(z)$ is self-reciprocal, so its roots come in pairs $\alpha$, $1/\bar\alpha$; on the circle the two members of a pair
have proportional distances to $z$, so keeping one member of each pair produces a square root of $T$. Once $T=|P|^2$, extremal problems
for $T$ become quadratic-form problems for the coefficients of $P$.

\writeup (a) If $T$ is constant, take $P=\sqrt{c_0}$. Otherwise reduce $n$ until $c_n\ne0$, and let $Q(z)=\sum_{k=-n}^nc_kz^{k+n}$, of degree $2n$ with $Q(0)=c_{-n}=\overline{c_n}\ne0$.
Since $c_{-k}=\overline{c_k}$, $z^{2n}\overline{Q(1/\bar z)}=Q(z)$ for $z\ne0$, so $\alpha\mapsto1/\bar\alpha$ permutes the roots of $Q$, preserving multiplicities. A root
$e^{i\theta_0}$ on the circle has even multiplicity: $T(\theta)=e^{-in\theta}Q(e^{i\theta})$ vanishes at $\theta_0$ to the order of that multiplicity and does not change
sign. So the roots can be listed as $\alpha_1,\dots,\alpha_n,1/\bar\alpha_1,\dots,1/\bar\alpha_n$ with $0<|\alpha_j|\le1$, and $Q(z)=c_n\prod_j(z-\alpha_j)(z-1/\bar\alpha_j)$. For $|z|=1$,
$|z-1/\bar\alpha|=\frac{|\bar\alpha z-1|}{|\alpha|}=\frac{|z|\,|\bar\alpha-\bar z|}{|\alpha|}=\frac{|z-\alpha|}{|\alpha|}$. Hence, as $T\ge0$,
\[
T(\theta)=|Q(e^{i\theta})|=\frac{|c_n|}{\prod_j|\alpha_j|}\prod_j|e^{i\theta}-\alpha_j|^2=|P(e^{i\theta})|^2,\qquad P(z)=\Bigl(\frac{|c_n|}{\prod_j|\alpha_j|}\Bigr)^{1/2}\prod_j(z-\alpha_j).
\]

(b) Write $P=\sum_{k=0}^np_kz^k$. By orthogonality, $c_0=\frac1{2\pi}\int|P(e^{i\theta})|^2\dd\theta=\sum|p_k|^2=1$, and by Cauchy--Schwarz
$T(0)=|\sum p_k|^2\le(n+1)\sum|p_k|^2=n+1$, with equality iff all $p_k$ are equal, that is $T=\frac1{n+1}|\sum_{k=0}^ne^{ik\theta}|^2$, the Fejér kernel $F_{n+1}$ of
Tome~4, Lemma~6.8.

(c) Now $c_k=\frac{a_k}2$ for $k\ge1$ and $c_0=1$. With $T=|P|^2$ as in (b), $c_1=\frac1{2\pi}\int|P|^2e^{-i\theta}\dd\theta=\sum_kp_{k+1}\overline{p_k}$, so with $u_k=|p_k|$,
\[
a_1=2\operatorname{Re}\sum_{k=0}^{n-1}p_{k+1}\overline{p_k}\le2\sum_{k=0}^{n-1}u_ku_{k+1}=u^{\mathsf T}Au,\qquad|u|=1,
\]
where $A$ is the adjacency matrix of the path with $n+1$ vertices. So $a_1\le\lambda_{\max}(A)$. The vectors $v^{(j)}_k=\sin\frac{(k+1)j\pi}{n+2}$, $0\le k\le n$, satisfy
$Av^{(j)}=2\cos\frac{j\pi}{n+2}v^{(j)}$ (by $\sin(x-y)+\sin(x+y)=2\sin x\cos y$, with $v_{-1}=v_{n+1}=0$), for $j=1,\dots,n+1$: these are $n+1$ distinct eigenvalues,
hence all of them, and the largest is $2\cos\frac\pi{n+2}$. For equality take $p_k=v^{(1)}_k/|v^{(1)}|>0$: then $T=|P|^2$ has real coefficients $c_k=c_{-k}$, so it is
a cosine polynomial $1+\sum a_k\cos k\theta$, and $a_1=2\sum p_kp_{k+1}=2\cos\frac\pi{n+2}$.

\tools{Tome~0, Chapter~3 (roots of polynomials) and Chapter~7 (symmetric matrices); Tome~4, Lemma~6.8 (the Fejér kernel).}
\leads Non-negative trigonometric polynomials are everywhere. Mertens's $3+4\cos\theta+\cos2\theta=2(1+\cos\theta)^2\ge0$ is the heart of the zero-free
region for $\zeta$ on the line $\operatorname{Re}s=1$ (Tome~5); Delsarte's linear programming bound for codes and the Cohn--Elkies bounds for sphere packings,
which Viazovska used in dimension $8$, optimise over non-negative or positive-definite functions; and engineers use the Fejér--Riesz
factorisation, under the name \emph{spectral factorisation}, to design filters.
\end{solution}

\begin{solution}{pb:bernstein}
\idea A trigonometric polynomial of degree $n$ has at most $2n$ zeros in a period. If $T$ climbed faster than the pure oscillation
$\cos(n(\theta-\beta))$ that has the same value at some point, their difference would have too many zeros.

\writeup (a) First suppose $\|T\|<1$, and suppose that $T'(\theta_0)^2+n^2T(\theta_0)^2>n^2$ for some $\theta_0$. Then $T'(\theta_0)\ne0$, as $|T(\theta_0)|<1$. Let $\gamma\in\left]0,\pi\right[$ with
$\cos\gamma=T(\theta_0)$, and choose $\beta$ with $n(\theta_0-\beta)=\pm\gamma$, the sign chosen so that $C(\theta)=\cos(n(\theta-\beta))$ has $C(\theta_0)=T(\theta_0)$ and $C'(\theta_0)$ of the
same sign as $T'(\theta_0)$; then $|C'(\theta_0)|=n\sin\gamma=n\sqrt{1-T(\theta_0)^2}<|T'(\theta_0)|$. The difference $D=C-T$ is a trigonometric polynomial of degree at most
$n$, and not $0$ since $C$ reaches $1$. At the $2n$ points $\theta_j=\beta+\frac{j\pi}n$ of a period, $C=(-1)^j$, so $D(\theta_j)$ has the sign of $(-1)^j$, and $D$ vanishes in
each of the $2n$ arcs $]\theta_j,\theta_{j+1}[$. Now look at the arc containing $\theta_0$, where $C$ is monotone; say $C$ decreases there, the other case
being symmetric. Then $D(\theta_j)>0>D(\theta_{j+1})$, and $T'(\theta_0)<C'(\theta_0)<0$, so $D'(\theta_0)>0$: $D$ is negative just before $\theta_0$ and positive just after. So $D$
has a zero in $]\theta_j,\theta_0[$, the zero $\theta_0$, and a zero in $]\theta_0,\theta_{j+1}[$: at least $2n+2$ zeros in a period, which is impossible. Hence
$T'^2+n^2T^2\le n^2$ when $\|T\|<1$, and by applying this to $(1-\varepsilon)T$ and letting $\varepsilon\to0$, also when $\|T\|\le1$. In particular $|T'|\le n$, and by
homogeneity $\|T'\|\le n\|T\|$. For complex $T$ and a given $\theta_0$, choose $\beta$ with $e^{i\beta}T'(\theta_0)=|T'(\theta_0)|$; the real polynomial $U=\operatorname{Re}(e^{i\beta}T)$ has
$\|U\|\le\|T\|$ and $U'(\theta_0)=|T'(\theta_0)|$, so $|T'(\theta_0)|\le n\|T\|$.

(b) $T(\theta)=P(\cos\theta)$ is a real trigonometric polynomial of degree $n$ with $\|T\|\le1$, and $T'(\theta)=-\sin\theta\,P'(\cos\theta)$. For $x=\cos\theta$ with
$0<\theta<\pi$, (a) gives $|P'(x)|\sin\theta\le n$, that is $|P'(x)|\le\frac n{\sqrt{1-x^2}}$.

(c) $T=\sin n\theta$ has $\|T\|=1$ and $T'(0)=n$. For the Chebyshev polynomial, $T_n(\cos\theta)=\cos n\theta$, so $T_n'(\cos\theta)=\frac{n\sin n\theta}{\sin\theta}$, and at $\theta=\frac\pi{2n}$,
$|T_n'(x)|=\frac n{\sqrt{1-x^2}}$.

\tools{Tome~0, Chapter~2 (trigonometric polynomials) and Chapter~3 (counting roots); Tome~4, Corollary~6.10 (Chebyshev polynomials).}
\leads Bernstein proved the inequality in 1912 to show that a function which trigonometric polynomials
of degree $n$ approximate within $n^{-\alpha}$ must be Hölder of order $\alpha$, for $0<\alpha<1$: the converse of Jackson's theorem. The sharper form
in (a) is the Bernstein--Szegő inequality. Near the ends of $[-1,1]$ the bound of (b) blows up, and the truth there is Markov's
inequality (1889): $|P'|\le n^2$ on $[-1,1]$, with equality for $T_n$ at $\pm1$. Problem~\ref{pb:bernsteinabs} uses (a) to build a function.
\end{solution}

\begin{solution}{pb:bernsteinabs}
\idea Hölder continuity controls the $L^2$ norm of differences, and Parseval converts this into $\ell^2$ control of each dyadic block of
coefficients. Cauchy--Schwarz turns $\ell^2$ into $\ell^1$ at the price of the square root of the block length; the books balance exactly
when $\alpha=\frac12$, and Rudin--Shapiro polynomials show that at $\alpha=\frac12$ they do not balance in our favour.

\writeup (a) $f$ is continuous, and $g(x)=f(x+h)-f(x-h)$ has coefficients $(e^{ikh}-e^{-ikh})c_k=2i\sin(kh)c_k$; apply Parseval (Tome~4, Theorem~6.11).

(b) The left side of (a) is at most $C^2(2h)^{2\alpha}$. For $j\ge0$ let $B_j=\{k:2^{j-1}<|k|\le2^j\}$, which has $2^j$ elements, and $h=\frac\pi{3\cdot2^j}$; for
$k\in B_j$, $\frac\pi6<|k|h\le\frac\pi3$, so $\sin^2(kh)\ge\frac14$ and
\[
\sum_{k\in B_j}|c_k|^2\le\sum_k4\sin^2(kh)|c_k|^2\le C^2\Bigl(\frac{2\pi}{3\cdot2^j}\Bigr)^{2\alpha}.
\]
By Cauchy--Schwarz,
\[
\sum_{k\in B_j}|c_k|\le2^{j/2}C\Bigl(\frac{2\pi}3\Bigr)^\alpha2^{-j\alpha}=C\Bigl(\frac{2\pi}3\Bigr)^\alpha2^{j(1/2-\alpha)},
\]
a geometric series in $j$ since $\alpha>\frac12$. So
$\sum_k|c_k|=|c_0|+\sum_j\sum_{B_j}|c_k|<\infty$. Then $\sum_kc_ke^{ikx}$ converges uniformly to a continuous function with the coefficients of $f$, which is $f$
(Tome~4, Corollary~6.10(b)).

(c) Let $g_j(x)=e^{i2^jx}P_j(e^{ix})$. By Problem~\ref{pb:rudin}, its frequencies are $2^j,\dots,2^{j+1}-1$, its coefficients are $\pm1$, and $\|g_j\|\le\sqrt{2^{j+1}}$. So
$f=\sum_j2^{-j}g_j$ converges uniformly, and integrating term by term, $c_k(f)=\pm2^{-j}$ for $2^j\le k<2^{j+1}$: $\sum_k|c_k(f)|=\sum_j2^{-j}2^j=\infty$. For the Hölder
bound, let $0<|h|\le1$ and $m\ge0$ with $2^{-m-1}<|h|\le2^{-m}$. As $g_j$ has degree less than $2^{j+1}$, Bernstein's inequality (Problem~\ref{pb:bernstein})
gives $\|g_j'\|\le2^{j+1}\sqrt{2^{j+1}}$. Hence
\[
\begin{aligned}|f(x+h)-f(x)|&\le\sum_{j\le m}2^{-j}|h|\,2^{3(j+1)/2}+\sum_{j>m}2^{-j}\cdot2\sqrt{2^{j+1}}\\&\le\frac{2^{3/2}|h|2^{m/2}+2^{3/2}2^{-(m+1)/2}}{1-2^{-1/2}}\le20|h|^{1/2},\end{aligned}
\]
since $2^{m/2}\le|h|^{-1/2}$ and $2^{-(m+1)/2}<|h|^{1/2}$. For $|h|>1$, $|f(x+h)-f(x)|\le2\|f\|\le2\|f\||h|^{1/2}$. So $f$ is Hölder of order $\frac12$.

\tools{Tome~4, Theorem~6.11 and Corollary~6.10; Problems~\ref{pb:rudin} and~\ref{pb:bernstein}.}
\leads The space of absolutely convergent Fourier series is an algebra, the Wiener algebra of Tome~7, whose invertible elements Wiener
characterised in 1932. Bernstein's theorem says that it contains every Hölder function of order $>\frac12$; (c), which is essentially
Rudin's motivation for his 1959 note, shows that order $\frac12$ is not enough. Hardy and Littlewood had shown this earlier with the series
$\sum_k\frac{e^{ick\log k}}ke^{ikx}$.
\end{solution}

\begin{solution}{pb:fantaussky}
\idea Write the polygon as a sum of regular polygons, one for each frequency $m$ (its discrete Fourier modes). Differences and areas
are both diagonal in this decomposition, and each mode is compared with its own constant; the extremal mode is $m=1$, the regular
polygon traversed once.

\writeup (a) Write $x_k=\sum_{m=0}^{n-1}c_m\omega^{mk}$ (DFT inversion, Tome~4, Theorem~7.3); $c_0=\frac1n\sum x_k=0$. Then $x_{k+1}-x_k=\sum_mc_m(\omega^m-1)\omega^{mk}$, and by
Parseval,
\[
\begin{aligned}\sum_k|x_{k+1}-x_k|^2&=n\sum_m|c_m|^2|\omega^m-1|^2=n\sum_{m=1}^{n-1}4\sin^2\frac{\pi m}n|c_m|^2\\&\ge4\sin^2\frac\pi n\cdot n\sum_m|c_m|^2=4\sin^2\frac\pi n\sum_k|x_k|^2,\end{aligned}
\]
since $\sin^2\frac{\pi m}n\ge\sin^2\frac\pi n$ for $1\le m\le n-1$, with equality only for $m=1$ and $m=n-1$. Equality holds iff $c_m=0$ for the other $m$, that is
$x_k=a\omega^k+b\omega^{-k}$.

(b) With $z_k=\sum_mc_m\omega^{mk}$, orthogonality gives $\sum_k\overline{z_k}z_{k+1}=n\sum_m|c_m|^2\omega^m$, so $2A=n\sum_m|c_m|^2\sin\frac{2\pi m}n$, while
$\sum_k|z_{k+1}-z_k|^2=n\sum_m|c_m|^2\,4\sin^2\frac{\pi m}n$. With $\theta_m=\frac{\pi m}n$,
\[
\sum_k|z_{k+1}-z_k|^2-4\tan\frac\pi n\cdot A=4n\sum_m|c_m|^2\sin\theta_m\Bigl(\sin\theta_m-\tan\frac\pi n\cos\theta_m\Bigr).
\]
The term $m=0$ vanishes. For $1\le m\le n-1$, $\sin\theta_m>0$; the bracket is positive if $\cos\theta_m\le0$, and equals $\cos\theta_m(\tan\theta_m-\tan\frac\pi n)\ge0$ if
$\cos\theta_m>0$, with equality only for $m=1$. So the difference is $\ge0$, with equality iff $z_k=c_0+c_1\omega^k$: a positively oriented regular $n$-gon
with centre $c_0$, or a point if $c_1=0$.

(c) If all sides have length $s$, then $L=ns$ and $\sum|z_{k+1}-z_k|^2=ns^2=\frac{L^2}n$, so (b) gives $A\le\frac{L^2}{4n\tan(\pi/n)}$, with equality for the regular
$n$-gon (orient the polygon positively, so that $A$ is its area).

\tools{Tome~4, Theorem~7.3 (DFT inversion and Parseval); Tome~0, Chapter~2 (roots of unity).}
\leads Ky Fan, Olga Taussky and John Todd (1955) proved (a) as a discrete analogue of Wirtinger's inequality, and (b) is the discrete
version of Hurwitz's Fourier proof of the isoperimetric inequality (Tome~4). The full polygonal isoperimetric inequality,
$L^2\ge4n\tan\frac\pi n\cdot A$ for every $n$-gon, needs a further argument (equalising the sides); the decomposition of a polygon into regular
modes is also the key to Problems~\ref{pb:midpoint} and~\ref{pb:napoleon}.
\end{solution}

\begin{solution}{pb:siegel}
\idea If $K$ avoids the lattice, the half-body $\frac12K$ packs, so the periodic function $\varphi$ is an indicator; its mean square equals its
mean, $2^{-n}\operatorname{vol}K$, but Bessel says the mean square is at least the square of the mean. This forces $2^{-n}\operatorname{vol}K\le1$, and the other Fourier
coefficients measure how far from tight the packing is.

\writeup (a) If $x+m=y+m'$ with $x,y\in\frac12K$ and $m\ne m'$, then $m-m'=y-x=\frac12(2y)+\frac12(-2x)$ with $2y\in K$ and $-2x\in K$ (symmetry), so
$m-m'\in K$ by convexity: a non-zero integer point of $K$, which is excluded.

(b) $K$ is bounded, so on $[0,1[^n$ only finitely many terms of $\varphi$ are not identically $0$, and by (a) at most one is non-zero at each point; so
$\varphi\in\{0,1\}$ and $\varphi^2=\varphi$. (A bounded convex set has a boundary of volume $0$, so these indicators are integrable.) Unfolding,
\[
\hat\varphi(k)=\sum_m\int_{[0,1[^n}1_{\frac12K}(x+m)e(-k\cdot x)\dd x=\sum_m\int_{[0,1[^n+m}1_{\frac12K}(y)e(-k\cdot y)\dd y=\int_{\frac12K}e(-k\cdot y)\dd y,
\]
since $e(k\cdot m)=1$ and the cubes $[0,1[^n+m$ tile $\R^n$. For $k=0$ this is $\operatorname{vol}\frac12K=2^{-n}\operatorname{vol}K$.

(c) The functions $e(k\cdot x)$ are orthonormal on the torus (the integrals factor into one-dimensional ones). For a finite $F\subset\Z^n$,
$0\le\int|\varphi-\sum_{k\in F}\hat\varphi(k)e(k\cdot x)|^2=\int\varphi^2-\sum_{k\in F}|\hat\varphi(k)|^2$. As $\int\varphi^2=\int\varphi=v$, with $v=2^{-n}\operatorname{vol}K$, we get
$v-v^2\ge\sum_{k\in F\setminus\{0\}}|\hat\varphi(k)|^2$ for every $F$, which is the inequality. In particular $v(1-v)\ge0$, so $\operatorname{vol}K\le2^n$; contrapositively, a symmetric
convex body of volume $>2^n$ contains a non-zero integer point.

(d) The open cube contains no non-zero integer point, and $v=1$. For $k\ne0$, some $k_j\ne0$, and $\hat\varphi(k)=\prod_j\int_{-1/2}^{1/2}e(-k_jx_j)\dd x_j=0$ because the
factor $j$ vanishes. Both sides are $0$: the translates of $\frac12K$ tile space.

\tools{Tome~4, Chapter~6 (orthogonality, Bessel's inequality) and Proposition~6.15 (Fourier analysis in several variables).}
\leads With Parseval in place of Bessel, the inequality becomes Siegel's identity, $v(1-v)=\sum_{k\ne0}|\hat\varphi(k)|^2$, which measures the defect $2^n-\operatorname{vol}K$
of a lattice-free body. Siegel published it in 1935, and ten years later his mean value theorem over the space of all lattices
opened a new chapter of the geometry of numbers (Tome~12). Minkowski's theorem itself, with its applications to sums of squares and
to class groups, is in Tome~1 and Tome~8, Appendix~I.
\end{solution}

\section*{Map of the chapter}

\begin{center}
\small
\adjustbox{max width=\linewidth}{\begin{tabular}{lll}
Problem & Topic & Tomes used\\\hline
\ref{pb:bilinear} & bilinear sums, $ab+cd=x$ in $\mathbb F_p$ & 0, 4\\
\ref{pb:kloosterman} & Kloosterman sums, fourth moment & 0, 1\\
\ref{pb:apcount} & progressions among quadratic residues & 1, 4\\
\ref{pb:dftspectrum} & eigenvalues of the DFT & 0, 4\\
\ref{pb:newman} & Newman's digit-sum bias & 0, 4\\
\ref{pb:erdosturan} & discrepancy of $k^2/p$ & 1, 4\\
\ref{pb:circlewalk} & a random walk on the circle & 3, 4\\
\ref{pb:kac} & lacunary sums look Gaussian & 0, 1, 3\\
\ref{pb:cantor} & the Cantor law has no density & 0, 3\\
\ref{pb:sincsum} & $\sum\operatorname{sinc}^kn$ versus $\int\operatorname{sinc}^k$ & 3, 4\\
\ref{pb:influence} & influences, majority & 0, 3\\
\ref{pb:granularity} & low degree, few variables & 0, 3\\
\ref{pb:bent} & bent functions & 2, 4\\
\ref{pb:majoritystable} & noise stability, Sheppard & 3, 4\\
\ref{pb:arrow} & Condorcet, Arrow, Kalai & 3\\
\ref{pb:bonami} & the Bonami lemma & 3\\
\ref{pb:hypercubemix} & cutoff on the hypercube & 3\\
\ref{pb:fejer} & non-negative trigonometric polynomials & 0, 4, 5\\
\ref{pb:bernstein} & Bernstein's inequality & 0, 4\\
\ref{pb:bernsteinabs} & absolutely convergent Fourier series & 4\\
\ref{pb:fantaussky} & polygons by the DFT & 0, 4\\
\ref{pb:siegel} & Minkowski's theorem by Fourier & 1, 4\\
\end{tabular}}
\end{center}

\section*{Sources and further reading}

\begin{itemize}
\item J.~Fourier, \emph{Théorie analytique de la chaleur}, Paris, 1822; English translation by A.~Freeman, \emph{The Analytical Theory of
Heat}, Cambridge, 1878 (the epigraph, from the Preliminary Discourse).
\item M.~J.~A.~N.~de Caritat, marquis de Condorcet, \emph{Essai sur l'application de l'analyse à la probabilité des décisions rendues à la
pluralité des voix}, Paris, 1785; K.~J.~Arrow, \emph{Social Choice and Individual Values}, Wiley, 1951; G.-Th.~Guilbaud, Les théories de
l'intérêt général et le problème logique de l'agrégation, \emph{Économie appliquée} 5 (1952), 501--584 (the limit $0.0877$, given ``by the
usual methods''; an English translation with commentary by B.~Monjardet appeared in the \emph{Journal Électronique d'Histoire des
Probabilités et de la Statistique} 4 (2008)).
\item G.~Kalai, A Fourier-theoretic perspective on the Condorcet paradox and Arrow's theorem, \emph{Advances in Applied Mathematics} 29
(2002), 412--426; E.~Friedgut, G.~Kalai and A.~Naor, Boolean functions whose Fourier transform is concentrated on the first two
levels, \emph{ibid.}, 427--437; E.~Mossel, A quantitative Arrow theorem, \emph{Probability Theory and Related Fields} 154 (2012), 49--88.
\item R.~O'Donnell, \emph{Analysis of Boolean Functions}, Cambridge University Press, 2014 (influences, noise stability, Kalai's formula,
the Bonami lemma, granularity; free online).
\item H.~D.~Kloosterman, On the representation of numbers in the form $ax^2+by^2+cz^2+dt^2$, \emph{Acta Mathematica} 49 (1926), 407--464;
A.~Weil, On some exponential sums, \emph{Proceedings of the National Academy of Sciences USA} 34 (1948), 204--207; N.~M.~Katz,
\emph{Gauss Sums, Kloosterman Sums, and Monodromy Groups}, Princeton University Press, 1988.
\item K.~F.~Roth, On certain sets of integers, \emph{Journal of the London Mathematical Society} 28 (1953), 104--109; F.~A.~Behrend, On sets
of integers which contain no three terms in arithmetical progression, \emph{Proceedings of the National Academy of Sciences USA} 32
(1946), 331--332; Z.~Kelley and R.~Meka, Strong bounds for 3-progressions, \emph{Proceedings of the 64th IEEE Symposium on Foundations of Computer Science} (2023).
\item J.~H.~McClellan and T.~W.~Parks, Eigenvalue and eigenvector decomposition of the discrete Fourier transform, \emph{IEEE
Transactions on Audio and Electroacoustics} 20 (1972), 66--74; I.~Schur, Über die Gaußschen Summen, \emph{Nachrichten der Gesellschaft
der Wissenschaften zu Göttingen} (1921), 147--153.
\item D.~J.~Newman, On the number of binary digits in a multiple of three, \emph{Proceedings of the AMS} 21 (1969), 719--721; J.~Coquet,
A summation formula related to the binary digits, \emph{Inventiones Mathematicae} 73 (1983), 107--115.
\item P.~Erdős and P.~Turán, On a problem in the theory of uniform distribution I, II, \emph{Indagationes Mathematicae} 10 (1948),
370--378 and 406--413; J.~D.~Vaaler, Some extremal functions in Fourier analysis, \emph{Bulletin of the AMS} 12 (1985), 183--216.
\item M.~Kac, On the distribution of values of sums of the type $\sum f(2^kt)$, \emph{Annals of Mathematics} 47 (1946), 33--49; R.~Salem and
A.~Zygmund, On lacunary trigonometric series, \emph{Proceedings of the National Academy of Sciences USA} 33 (1947), 333--338; M.~Kac,
Probability methods in some problems of analysis and number theory, \emph{Bulletin of the AMS} 55 (1949), 641--665 (the
Erdős--Fortet example).
\item P.~Erdős, On a family of symmetric Bernoulli convolutions, \emph{American Journal of Mathematics} 61 (1939), 974--976; B.~Solomyak,
On the random series $\sum\pm\lambda^n$ (an Erdős problem), \emph{Annals of Mathematics} 142 (1995), 611--625.
\item R.~Baillie, D.~Borwein and J.~M.~Borwein, Surprising sinc sums and integrals, \emph{American Mathematical Monthly} 115 (2008),
888--901.
\item J.~Kahn, G.~Kalai and N.~Linial, The influence of variables on Boolean functions, \emph{Proceedings of FOCS 1988}, 68--80;
N.~Nisan and M.~Szegedy, On the degree of Boolean functions as real polynomials, \emph{Computational Complexity} 4 (1994), 301--313;
J.~Chiarelli, P.~Hatami and M.~Saks, An asymptotically tight bound on the number of relevant variables in a bounded degree Boolean
function, \emph{Combinatorica} 40 (2020), 237--244.
\item O.~S.~Rothaus, On ``bent'' functions, \emph{Journal of Combinatorial Theory A} 20 (1976), 300--305.
\item W.~F.~Sheppard, On the application of the theory of error to cases of normal distribution and normal correlation,
\emph{Philosophical Transactions of the Royal Society A} 192 (1899), 101--167; E.~Mossel, R.~O'Donnell and K.~Oleszkiewicz, Noise
stability of functions with low influences: invariance and optimality, \emph{Annals of Mathematics} 171 (2010), 295--341.
\item A.~Bonami, Étude des coefficients de Fourier des fonctions de $L^p(G)$, \emph{Annales de l'Institut Fourier} 20 (1970), 335--402;
W.~Beckner, Inequalities in Fourier analysis, \emph{Annals of Mathematics} 102 (1975), 159--182.
\item D.~Aldous and P.~Diaconis, Shuffling cards and stopping times, \emph{American Mathematical Monthly} 93 (1986), 333--348;
P.~Diaconis, \emph{Group Representations in Probability and Statistics}, IMS, 1988 (the hypercube walk, Chapter~3).
\item L.~Fejér, Über trigonometrische Polynome, \emph{Journal für die reine und angewandte Mathematik} 146 (1916), 53--82; E.~Egerváry
and O.~Szász, Einige Extremalprobleme im Bereiche der trigonometrischen Polynome, \emph{Mathematische Zeitschrift} 27 (1928), 641--652.
\item S.~Bernstein, Sur l'ordre de la meilleure approximation des fonctions continues par des polynômes de degré donné, \emph{Mémoires
de l'Académie royale de Belgique} (2) 4 (1912); P.~Borwein and T.~Erdélyi, \emph{Polynomials and Polynomial Inequalities}, Springer, 1995.
\item W.~Rudin, Some theorems on Fourier coefficients, \emph{Proceedings of the AMS} 10 (1959), 855--859; A.~Zygmund,
\emph{Trigonometric Series}, Cambridge University Press, 1959 (Bernstein's theorem, Chapter~VI).
\item K.~Fan, O.~Taussky and J.~Todd, Discrete analogs of inequalities of Wirtinger, \emph{Monatshefte für Mathematik} 59 (1955), 73--90.
\item C.~L.~Siegel, Über Gitterpunkte in convexen Körpern und ein damit zusammenhängendes Extremalproblem, \emph{Acta Mathematica} 65
(1935), 307--323.
\end{itemize}
